% !TEX program = xelatex
%% presentation_script.tex : 발표 대본(1부)과 발표자 배경 지식(2부). 2026-10-06 연구실 보고용.
%% 컴파일: ./build.sh (xelatex 2회, nice -n 10). 출력 deck/render/발표대본_연구결과보고_2026-10-06.pdf
%% 글꼴: Pretendard(본문·제목, paper/manuscript/ko/main.tex 의 제목 글꼴과 같음), DejaVu Sans Mono(경로).
%% 구성: 1부 = part1_script.tex(슬라이드별 대본), 2부 = part2_background.tex(배경 지식).
\documentclass[11pt,a4paper]{article}

% ---------- 레이아웃 ----------
\usepackage[a4paper,top=24mm,bottom=26mm,left=24mm,right=24mm,footskip=13mm]{geometry}
\usepackage{amsmath}
\usepackage{amssymb}
\usepackage{fontspec}
\usepackage{kotex}   % xetexko: 한글·라틴 사이에 자동 공백을 넣지 않고 입력한 공백을 지킨다
\usepackage{xcolor}
\usepackage{array}
\usepackage{booktabs}
\usepackage{tabularx}
\usepackage{longtable}
\usepackage{enumitem}
\usepackage{titlesec}
\usepackage{fancyhdr}
\usepackage[most]{tcolorbox}
\usepackage[hidelinks]{hyperref}

% ---------- 글꼴 ----------
\setmainfont{Pretendard}[
  UprightFont={Pretendard}, BoldFont={Pretendard Bold},
  ItalicFont={Pretendard}, ItalicFeatures={FakeSlant=0.15},
  BoldItalicFont={Pretendard Bold}, BoldItalicFeatures={FakeSlant=0.15}]
\setsansfont{Pretendard}[BoldFont={Pretendard Bold}]
\setmonofont{DejaVu Sans Mono}[Scale=0.86]
\setmainhangulfont{Pretendard}[
  UprightFont={Pretendard}, BoldFont={Pretendard Bold},
  ItalicFont={Pretendard}, ItalicFeatures={FakeSlant=0.15},
  BoldItalicFont={Pretendard Bold}, BoldItalicFeatures={FakeSlant=0.15}]
\setsanshangulfont{Pretendard}[BoldFont={Pretendard Bold}]
\setmonohangulfont{Pretendard}[Scale=0.9]
\newfontfamily\semibold{Pretendard SemiBold}
\newhangulfontfamily\cjksemibold{Pretendard SemiBold}

\linespread{1.32}
\setlength{\parindent}{0pt}
\setlength{\parskip}{0.55em}
\setlength{\emergencystretch}{2em}
\raggedbottom

% ---------- 색 ----------
\definecolor{accent}{HTML}{C8651B}
\definecolor{rulegray}{HTML}{9A9A9A}
\definecolor{textgray}{HTML}{555555}
\definecolor{boxbg}{HTML}{F4F4F2}
\definecolor{boxframe}{HTML}{BDBDBD}

% ---------- 절 서식 ----------
\setcounter{secnumdepth}{2}
\titleformat{\section}{\semibold\cjksemibold\Large}{\thesection}{0.6em}{}
\titleformat{\subsection}{\semibold\cjksemibold\large}{\thesubsection}{0.6em}{}
\titleformat{\subsubsection}{\semibold\cjksemibold\normalsize}{\thesubsubsection}{0.6em}{}
\titlespacing*{\section}{0pt}{2.4ex plus 0.6ex}{1.0ex}
\titlespacing*{\subsection}{0pt}{1.8ex plus 0.4ex}{0.6ex}
\titlespacing*{\subsubsection}{0pt}{1.4ex plus 0.3ex}{0.4ex}

% ---------- 머리글·쪽 번호 ----------
\pagestyle{fancy}
\fancyhf{}
\renewcommand{\headrulewidth}{0pt}
\fancyhead[L]{\small\color{textgray}\leftmark}
\fancyhead[R]{\small\color{textgray}연구 결과 종합 보고 · 2026-10-06}
\fancyfoot[C]{\small\thepage}
\renewcommand{\sectionmark}[1]{\markboth{#1}{}}
\fancypagestyle{plain}{\fancyhf{}\fancyfoot[C]{\small\thepage}\renewcommand{\headrulewidth}{0pt}}

% ---------- 1부 대본용 명령 ----------
\newcommand{\needlines}[1]{\par\vskip 0pt plus #1\penalty-100\vskip 0pt plus -#1\relax}
\newcommand{\slideheading}[3]{%
  \par\needlines{7\baselineskip}\vspace{1.6em plus 0.4em}\penalty-600
  \noindent{\semibold\cjksemibold\color{accent}\large 슬라이드 #1 · #2}\hfill\mbox{\ \small\color{textgray}#3분}\par
  \vspace{0.2em}\noindent{\color{rulegray}\rule{\textwidth}{0.4pt}}\par
  \vspace{0.2em}\markboth{1부 발표 대본}{}}
\newcommand{\pointat}[1]{\par\vspace{-0.1em}\noindent{\small\color{textgray}▶ 이 쪽에서 가리킬 것: #1}\par}
\newtcolorbox{optbox}{enhanced, breakable, colback=boxbg, colframe=boxframe, boxrule=0.4pt, arc=1pt,
  left=6pt, right=6pt, top=4pt, bottom=4pt, before skip=5pt, after skip=7pt,
  fontupper=\small, title={시간이 있으면}, fonttitle=\small\semibold\cjksemibold, coltitle=black,
  colbacktitle=boxbg, attach boxed title to top left={xshift=6pt, yshift=-2.2mm},
  boxed title style={boxrule=0pt, arc=0pt, colback=boxbg, left=2pt, right=2pt, top=0pt, bottom=0pt}}

% ---------- 2부용 표 서식 ----------
\newcolumntype{Y}{>{\raggedright\arraybackslash}X}
\newcommand{\tabsmall}{\small\setlength{\tabcolsep}{5pt}\renewcommand{\arraystretch}{1.18}}
\newcommand{\pth}[1]{{\ttfamily\small #1}}
\setlength{\LTleft}{0pt}\setlength{\LTright}{\fill}\setlength{\LTpre}{0.6em}\setlength{\LTpost}{0.8em}

% ---------- 부 표지 ----------
\newcommand{\partpage}[3]{%
  \clearpage\thispagestyle{empty}
  \vspace*{0.30\textheight}
  {\noindent\color{accent}\semibold\cjksemibold\Huge #1\par}
  \vspace{1.2em}
  {\noindent\semibold\cjksemibold\Large #2\par}
  \vspace{1.2em}
  {\noindent\color{textgray}#3\par}
  \clearpage}

\begin{document}

% ---------- 표지 ----------
\thispagestyle{empty}
\vspace*{0.16\textheight}
{\noindent\color{accent}\semibold\cjksemibold\huge 발표 대본과 발표자 배경 지식\par}
\vspace{1.0em}
{\noindent\semibold\cjksemibold\Large 희소 관측 영구동토 지역의 활동층 두께 예측을 위한\\물리 기반 기계학습 워크플로\par}
\vspace{1.0em}
{\noindent\large 연구 결과 종합 보고 · 2026년 10월 6일 · 연구실 세미나(20–30분)\par}
\vspace{2.2em}
{\noindent 발표 백승원 · 서울대학교 에너지시스템공학부 석박통합과정\par}
\vspace{2.6em}
{\noindent\color{textgray}\small
덱: \pth{deck/render/permafrost\_paper\_report.pdf}(재배열 최종판: 표지, 본편 1–36쪽, 마무리 37쪽, 참고문헌, 부록. 쪽 순서는 덱 스펙 slides 순서).\\
1부는 본편 슬라이드마다 말할 내용(계획 시간 합 25분)이고, 2부는 말하지 않는 배경 지식(경과, 코드, 자료, 방법, 결과 해석, 예상 질문, 수치 요약)이다.\\
수치는 덱 스펙(\pth{deck/deck\_spec\_paper\_report.json}), 최종 묶음 계획서 8절과 추가 등록 두 건, \pth{paper/claims/*/README.md}에서 옮겼고 새로 계산한 값은 없다.\par}
\vfill
{\noindent\color{textgray}\small 작성 2026-10-05. 2026-10-06 재배열 덱에 맞춰 1부를 다시 썼다(원문 \pth{deck/script/gen\_part1.py}, 같은 원문을 덱 스펙 notes.script 에 넣음). 저장소에는 커밋하지 않았다.\par}
\clearpage

% ---------- 차례 ----------
\thispagestyle{plain}
{\noindent\semibold\cjksemibold\Large 차례\par}
\vspace{0.6em}
{\noindent\semibold\cjksemibold 1부 발표 대본\par}
\begin{itemize}[leftmargin=1.4em, itemsep=0.1em]
  \item 표지, 1부 연구 배경(슬라이드 1–3), 2부 연구 목적(4–5): 3.4분
  \item 3부 연구 방법(6–15): 6.4분
  \item 4부 연구 결과: 라벨이 많은 지역(16–20) 3.6분, 새 지역(21–25) 3.1분, 단계적 예측 워크플로(26–30) 3.7분
  \item 5부 ALT 지도 비교(31–32): 0.9분
  \item 6부 결론(33–36), 마무리(37): 3.9분
\end{itemize}
{\noindent\semibold\cjksemibold 2부 발표자 배경 지식\par}
\begin{enumerate}[leftmargin=1.8em, itemsep=0.1em]
  \item 연구 수행 경과
  \item 코드 구성
  \item 자료
  \item 방법 상세
  \item 결과가 왜 이렇게 나왔는가
  \item 발표에서 말하지 않은 것
  \item 예상 질문과 답
  \item 수치 요약표
  \item 출처를 찾지 못해 뺀 수치
\end{enumerate}

\vspace{1em}
{\noindent\semibold\cjksemibold 대본 읽는 법\par}
\begin{itemize}[leftmargin=1.4em, itemsep=0.1em]
  \item 머리글의 분 수는 계획 시간이다. 합이 25.3분이며 질의응답은 포함하지 않는다.
  \item 회색 상자 '시간이 있으면'은 선택 문장이다. 시간이 부족하면 건너뛴다.
  \item '이 쪽에서 가리킬 것'은 레이저 포인터로 짚을 요소다.
  \item 판정어(오차 감소, 오차 증가, 동등, 미결정)와 근거 종류(사전 등록, 사후 분석, 사후 서술)는 덱과 같은 뜻으로 쓴다.
\end{itemize}

% ---------- 1부 ----------
\partpage{1부}{발표 대본}{재배열 최종 덱 · 본편 36쪽과 마무리 · 계획 시간 25분 · 정성적 설명 중심}
\markboth{1부 발표 대본}{}
% part1_script.tex : 발표 대본(슬라이드별). presentation_script.tex 에서 \input 한다.
% 본문 수치는 deck/deck_spec_paper_report.json(notes.script, numbers)과
% docs/EXPERIMENT_PLAN_FINAL_BATCH_2026-10-04.md 8절, 두 추가 등록 문서, paper/claims/*/README.md 에서 옮겼다.
% 분량 기준: 약 300자/분(공백 제외). 선택 상자는 분량에 넣지 않는다. 글자 수는 count_part1.py 로 센다.

\slideheading{표지}{희소 관측 영구동토 지역의 활동층 두께 예측을 위한 물리 기반 기계학습의 라벨 수별 워크플로}{0.3}

안녕하십니까, 백승원입니다. 현장 관측값, 즉 라벨이 적은 영구동토 지역에서 활동층 두께를 예측할 때, 라벨 수에 따라 물리식과 기계학습을 어떻게 결합하고 평가할지를 정리한 결과를 보고드리겠습니다. 배경, 목적, 방법, 결과, 지도, 결론 순서입니다.

\pointat{바탕 지도(이 방법으로 만든 알래스카 1 km 지도), 제목의 ‘라벨 수별 워크플로’}

\slideheading{1}{활동층 두께와 관측 공백}{1.0}

먼저 대상입니다. 왼쪽 a처럼 영구동토 위 지표층은 여름에 녹았다가 겨울에 다시 업니다. 그 계절 최대 융해 깊이가 활동층 두께, ALT이고 단위는 cm입니다. 기반시설과 탄소·수문 과정에 영향을 주므로 넓은 영역의 지도가 필요합니다. 그런데 가운데 b처럼 관측은 점 관측망의 수백 지점에 몰려 있습니다. 직접 라벨은 1~km 셀로 묶어 17,467개이고 학습 원천 셀의 78에서 94\%가 알래스카입니다. 오른쪽 c의 네 지도는 같은 알래스카인데 영역 평균이 56에서 105~cm로 제품마다 다릅니다. 즉 대부분의 영구동토는 라벨이 거의 없는 새 지역이고, 어느 지도를 믿을지도 정해져 있지 않습니다.

\pointat{a 의 ‘융해 전선’과 ALT 화살표, b 의 출처 표, c 의 평균값 네 개}

\begin{optbox}
b 아래 수치처럼 지역별 Stefan 계수의 기하 평균은 레나델타 1.31에서 티베트 고원 11.03까지 다릅니다. 같은 열량에도 녹는 깊이가 지역마다 다르다는 뜻이고, 계수 재보정이 필요한 이유입니다.
\end{optbox}

\slideheading{2}{선행 연구와 남은 질문}{0.6}

ALT 기계학습 선행 연구 10편을 평가 방식으로 비교한 표입니다. 위 다섯 편은 학습 지역 안 무작위 교차검증, 아래 네 편은 거리 블록이나 지점 제외의 공간 교차검증입니다. 공통점이 셋입니다. 대상 지역을 통째로 빼고 평가한 연구가 없고, 라벨 수는 0개 또는 전량의 한 조건이며, 같은 라벨로 재보정한 물리 기준선이 없습니다. 맨 아래 주황색 행이 이 연구입니다. 남은 질문은 새 지역에서 라벨 수에 따라 무엇이 통하는가입니다.

\pointat{‘평가 유형’ 열, 맨 아래 주황색 행}

\begin{optbox}
조사 범위는 2019년부터 2026년까지의 영문 ALT 기계학습 논문이며 본문과 요약을 확인한 수준입니다. 이 범위에서 찾지 못했다는 뜻이지 ‘최초’라는 주장은 아닙니다.
\end{optbox}

\slideheading{3}{기존 연구의 한계와 이 연구의 개선}{0.7}

세 행이 평가 지역, 라벨 수, 물리 기준선입니다. 평가 지역은 지역 홀드아웃으로 바꿉니다. 대상 지역과 경계 100~km 완충을 학습에서 빼므로 학습에 없는 지역의 오차를 추정합니다. 라벨 수는 0개부터 전량까지의 곡선으로 평가해 몇 개부터 기계학습을 더할지 판단합니다. 물리 기준선은 같은 라벨로 재보정한 Stefan 식입니다. 물리식도 같은 라벨을 써야 재보정 몫과 기계학습 몫이 분리되기 때문입니다. 왼쪽 위 그래프처럼 같은 직접 ML이 무작위 분할에서는 15~cm 남짓, 지역 홀드아웃에서는 29~cm 가까이로 커집니다.

\pointat{왼쪽 위 그래프의 두 선, 가운데 열의 주황색 세 항목}

\slideheading{4}{연구 목적과 문제 정의}{0.7}

왼쪽 지도에서 원천 지역은 학습에 쓰는 지역, 대상 지역은 예측하려는 새 지역이고, 대상 지역과 100~km 완충 안의 셀은 학습에서 뺍니다. 원천 자료와 Stefan 식이 있고 대상 라벨이 $n$개일 때 무엇을 학습하고 무엇으로 평가하는가, 이것이 문제입니다. 오른쪽은 비교할 네 방법입니다. 원천 계수 Stefan은 대상 라벨을 쓰지 않고, 재보정 Stefan은 계수 $E$만 고칩니다. 물리 잔차 결합은 재보정 Stefan을 기본 답으로 두고 잔차를 학습하며, 직접 ML은 물리 결합 없이 배웁니다. 모두 같은 라벨과 같은 채점 블록을 씁니다.

\pointat{삽도의 100 km 완충 점선, 오른쪽의 라벨 수 축}

\slideheading{5}{연구 질문}{0.5}

질문은 라벨 수 구간별로 넷입니다. 라벨이 없을 때 물리 정보가 전이 오차를 줄이는가. 라벨 몇 개로 무엇을 먼저 고치는가. 라벨이 늘면 방법을 어떻게 고르는가. 어떻게 평가하고 지도는 무엇을 담는가. 오른쪽 열은 답을 한 줄로 미리 적은 것이고 근거는 결과에서 하나씩 보이겠습니다.

\pointat{오른쪽 ‘요지’ 열}

\slideheading{6}{용어}{0.6}

용어를 먼저 정합니다. 라벨은 현장에서 잰 ALT 관측값으로 1~km 셀 단위입니다. 원천 지역은 학습 지역, 대상 지역은 새 지역입니다. 원천 계수는 원천 라벨로 정한 Stefan 계수, 재보정 계수는 대상 라벨로 다시 맞춘 계수입니다. 물리 잔차 결합은 물리식이 남긴 오차를 기계학습이 보정하는 이 연구의 제안이고, 직접 ML은 물리 정보 없이 배웁니다. 판정은 셀 가중과 블록 등가중의 두 95\% 신뢰구간으로 하고, 차이가 ±0.5~cm 안이면 같다고 보는 것이 동등 범위입니다.

\pointat{‘물리 잔차 결합’ 행, ‘두 가중 신뢰구간’과 ‘동등 범위’ 행}

\slideheading{7}{사용 자료}{0.8}

자료는 라벨, 공변량, 물리 입력의 셋입니다. 가운데 표에서 주 4지역은 레나델타 점 관측 3,037행, 캐나다 750행, 러시아 서부와 동부의 CALM 지점 평균 31행과 30행입니다. 참조 지역 알래스카는 13,606행이지만 1~km 위치 343곳, 0.5° 블록 74개에 몰려 있습니다. 새 지역 러시아 중부와 티베트 고원은 57행과 132행입니다. 라벨 수 $n$은 행 수이고 행의 뜻이 지역마다 다릅니다. 공변량은 지형 6종, 기후 8종, 토양 9종, 위성 ALT 2종의 25종이며, 융해 도일은 ERA5-Land에서 계산합니다. 오른쪽 c처럼 모두 1~km 셀에 정렬하지만 정보 해상도는 토양 약 5~km, 기후 0.1°에 묶입니다.

\pointat{표의 알래스카 행(13,606 / 343 / 74), c 의 ‘정보 해상도’ 띠}

\begin{optbox}
라벨 연도는 1990년부터 2024년까지이고 예측은 2015년부터 2020년 기후값에 대한 정적 예측입니다. 연별 변동은 지금 공변량으로 예측되지 않아 다년 평균 규약을 씁니다.
\end{optbox}

\slideheading{8}{전체 연구 흐름}{0.6}

여섯 단계입니다. 자료를 1~km 셀로 정리하고, Stefan 식으로 원천 계수와 재보정 계수의 두 기준선을 만듭니다. 셋째, 기계학습을 물리 정보와 결합하는 방식들을 같은 학습기 CatBoost로 비교합니다. 넷째, 모든 방법을 지역 홀드아웃과 라벨 수 격자에서 같은 채점 블록으로 평가하며 판정 규칙은 실행 전에 등록했습니다. 다섯째, 결과를 라벨 수별 워크플로로 묶고, 여섯째, 알래스카 1~km 지도와 기존 제품 비교로 마무리합니다.

\pointat{단계 ①–⑥ 띠, 아래 주황색 되돌림 화살표(라벨을 더 모으면 앞 단계로)}

\slideheading{9}{Stefan 물리 모델의 뜻}{0.9}

Stefan 식은 여름에 지면으로 들어온 융해 열량으로 녹는 깊이를 구하는 열수지 해입니다. 열량은 영상 기온의 합인 융해 도일 TDD로 대신하고, 열전도도, 토양 수분, 밀도, 융해 잠열은 계수 $E$ 하나에 묶입니다. 그래서 ALT는 $E$ 곱하기 루트 TDD입니다. 왼쪽 아래는 러시아 서부 라벨 셀 하나인데, TDD 1,537도·일에 실측 118~cm를 원천 계수는 62~cm로 예측합니다. $E$가 지역마다 다르기 때문입니다. 오른쪽은 러시아 서부 라벨 31개에서 10개를 뽑아 계수를 다시 맞춘 예시로, 원천 계수 1.59가 2.11로 바뀝니다. 이것이 다음 쪽의 재보정입니다.

\pointat{가운데 식의 $E=\sqrt{2k/(\rho w L)}$, 오른쪽의 회색 선(1.59)과 파란 선(2.11)}

\begin{optbox}
$E$의 단위는 cm·(°C·일)$^{-1/2}$이고 주 지역의 원천 계수는 1.59에서 1.62로 거의 같습니다. 원천의 대부분이 알래스카 셀이기 때문입니다.
\end{optbox}

\slideheading{10}{Stefan 물리 모델과 계수 재보정}{0.8}

재보정 계수 $E_n$은 대상 라벨 $n$개의 최소제곱 계수 $E_{\mathrm{ls}}$와 원천 계수 $E_0$의 가중 평균이고 수축 강도 $\kappa$는 10입니다. 왜 섞느냐면, 라벨 몇 개로 계수를 통째로 바꾸면 오차가 오히려 커질 수 있기 때문입니다. 원천 계수에 라벨 10개의 무게를 주면 라벨 10개에서 두 계수의 중간이 되고 라벨이 늘수록 대상 계수에 가까워집니다. $\kappa$는 앞선 실험에서 정했고 3과 30은 민감도입니다. 그림은 레나델타 실제 라벨에서 10개를 뽑은 예시이며, 재보정 직선과 라벨의 차이인 잔차를 다음 쪽의 기계학습이 배웁니다.

\pointat{아래 식 $E_n=(nE_{\mathrm{ls}}+\kappa E_0)/(n+\kappa)$, 회색 원천 계수 선과 파란 재보정 선, 주황색 ‘잔차’}

\begin{optbox}
이 예시의 최소제곱 계수는 1.352, 재보정 계수는 1.486, 원천 계수는 1.619입니다. 라벨 10개이므로 재보정 계수가 정확히 중간입니다. 판정과 무관한 예시입니다.
\end{optbox}

\slideheading{11}{비교 방법과 모델 구조}{1.0}

비교한 방법은 물리 정보가 들어가는 위치로 나눈 다섯 구조입니다. 위 두 줄이 원천 계수와 재보정 Stefan이고, 주황색이 제안하는 물리 잔차 결합입니다. 재보정 Stefan 예측 $a$를 앵커로 두고 학습기 $g$는 라벨과 앵커의 차이인 잔차만 배웁니다. 예측은 $a$ 더하기 $\lambda$ 곱하기 $g(X)$이고 잔차 가중 $\lambda$는 0.25, 0.5, 1.0 가운데 고르며 전이 판정 기준은 0.25입니다. 배운 것이 없으면 예측이 물리식으로 돌아가는 구조입니다. 물리 입력 ML은 물리 예측을 입력 열로 넣는데 학습기가 그 열을 무시할 수 있습니다. 물리 유사라벨 증강은 Stefan 값을 학습 라벨로 더하고, 직접 ML은 물리 정보를 쓰지 않습니다. 모두 같은 CatBoost로 비교했고, 라벨 0개에서는 학습기 10종으로 결론이 학습기에 묶이지 않는지 확인했습니다.

\pointat{주황색 물리 잔차 결합 줄($\hat{y}=a+\lambda g(X)$), 오른쪽의 $\lambda$ 설명}

\begin{optbox}
CatBoost 설정은 반복 200회, 학습률 0.05, 깊이 3, seed 2개입니다. 학습기 10종은 CatBoost 3설정, 랜덤 포레스트, TabPFN v2, TabICL v2, 신경망 4종입니다. 공변량의 CCI ALT 2종은 같은 재분석 자료로 구동한 모형 출력이라 관측으로 취급하지 않습니다.
\end{optbox}

\slideheading{12}{물리 유사라벨 증강 대조 실험 설계}{0.6}

증강의 이득이 물리 정보에서 오는지, 학습 행이 늘어난 효과인지를 가르는 대조 실험입니다. 유사라벨의 자리와 수, 학습기, 채점 블록을 같게 두고 값만 바꾼 위약 넷을 만들었습니다. 순서를 섞은 유사라벨, 풀 평균 상수, 원천 평균 상수, TDD에 선형인 융해 지수입니다. 유사라벨은 라벨 블록 셀에 원천 행당 10개를 넣습니다. 오른쪽은 물리 유사라벨에서 위약을 뺀 오차 변화이고 음수가 물리 유사라벨이 나은 쪽입니다. 해석은 결과에서 하겠습니다.

\pointat{‘유사라벨 조건’ 다섯 줄, 오른쪽 그래프의 방향}

\slideheading{13}{평가 설계}{1.0}

핵심이므로 자세히 말씀드립니다. a는 지역 홀드아웃입니다. 대상 지역 셀과 100~km 완충을 학습 원천에서 뺍니다. b는 블록 절반 분할입니다. 0.5° 블록을 무작위로 절반씩 나누어 라벨은 노란 블록에서만 뽑고 채점은 초록 블록에서만 하며, 분할은 다섯 번입니다. c는 라벨 수 격자로 0, 3, 10, 40, 160, 320, 1000개와 전량이고 추출도 다섯 번씩입니다. d는 판정 규칙입니다. 두 가중의 신뢰구간이 모두 0 아래면 오차 감소, 모두 0 위면 오차 증가, 네 끝값이 ±0.5~cm 안이면 동등, 그 밖은 미결정입니다. 미결정은 차이를 확인하지 못했다는 뜻입니다. 맨 아래 줄처럼 가설과 판정 규칙을 커밋한 뒤 고정 seed로 실행하고, 재현 관문을 통과한 뒤에 봉인한 표를 열어 그 시각을 기록했습니다.

\pointat{b 의 노란 블록과 초록 블록, d 의 네 판정 그림, 맨 아래 ‘등록 → 실행 → 재현 관문 → 열람’}

\begin{optbox}
러시아 서부와 동부는 라벨 블록 셀이 15개 안팎이라 40개 이상은 레나델타와 캐나다 2지역 평균으로 봅니다. ±0.5~cm는 주 4지역 원천 계수 Stefan RMSE 21.7에서 43.0~cm의 약 1에서 2\%이며 결과 전에 고정한 규약입니다.
\end{optbox}

\slideheading{14}{검증 설계에 따른 오차}{0.7}

왜 지역 홀드아웃인지를 수치로 보입니다. 가로축은 시험 셀에서 가장 가까운 학습 셀까지의 거리, 세로축은 RMSE입니다. 보라색 파선인 직접 ML의 3지역 평균은 무작위 셀 분할 14.8~cm에서 지역 홀드아웃 28.7~cm로 거의 두 배가 되지만, 파란 실선인 재보정 Stefan은 20.5에서 22.2~cm로 거의 변하지 않습니다. kNNDM은 시험 셀과 학습 셀의 거리 분포를 지도 예측 상황에 맞춘 분할입니다. 무작위 분할은 학습 셀이 시험 셀 옆에 있어 새 지역 성능을 과대평가합니다.

\pointat{가로축의 거리 눈금, 두 선이 교차하는 지점}

\begin{optbox}
무작위에서 kNNDM으로 갈 때 직접 ML은 알래스카, 레나델타, 캐나다에서 5.8, 8.2, 14.2~cm, 재보정 Stefan은 0.3, 0.5, 4.3~cm 늘었습니다. 확률 표본이라면 무작위 검증도 타당하다는 반론이 있으나 이 연구의 라벨은 관측망에 몰린 자료라 확률 표본이 아닙니다.
\end{optbox}

\slideheading{15}{라벨 수별 워크플로 알고리즘}{0.5}

결과를 새 지역에 적용하는 절차입니다. 라벨이 없으면 원천 계수 Stefan을 쓰고 학습 범위 밖 셀을 표시합니다. 3개 이상이면 계수를 재보정하고 낮은 가중의 잔차를 더하며, 10개에서는 원천 계수 Stefan의 평균 절대 오차로 편향을 진단합니다. 40개 이상이면 라벨 블록 5겹 교차검증으로 후보 7개 가운데 하나를 고릅니다. 라벨 배치 효과는 지역에 따라 달라 사전 지정 절차는 무작위 배치로 확정되었습니다.

\pointat{흐름도의 네 상자, ‘교차검증 후보 7개’ 줄, ‘라벨 배치’ 줄}

\slideheading{16}{라벨 수 구간별 결과 요지}{0.5}

라벨 0개에서는 30개 대상 가운데 2~cm 넘게 나빠진 대상이 물리 유사라벨 증강 2개, 직접 ML 18개였습니다. 사후 서술입니다. 3에서 10개에서는 이득의 대부분이 계수 재보정이었습니다. 40에서 160개에서는 교차검증 선정이 고정 레시피에 뒤지지 않았고, 수백 개 이상인 알래스카에서는 물리 잔차 결합이 0.54~cm 작았으며 이득은 기후 격자 단위 보정이었습니다. 분산 배치 효과는 지역에 따라 달랐습니다.

\pointat{표의 다섯 행}

\slideheading{17}{물리 유사라벨의 정보 출처}{0.5}

위약 대조의 결과입니다. 라벨 0개에서 순서를 섞은 유사라벨보다 오차가 1.63~cm 작았고, 두 신뢰구간 모두 0 아래이며 보정 뒤 $p$는 0.001입니다. 다만 맨 아래 선형 융해 지수와의 차이는 0.48~cm로 0.5~cm 미만이고 라벨 10개에서는 동등입니다. 이득은 셀마다 다른 물리 값에서 오지만 대부분은 융해 도일에 선형인 정보로도 얻어집니다. 사전 등록 판정입니다.

\pointat{왼쪽이 물리 유사라벨이 나은 쪽, 회색 띠 ±0.5 cm. 첫 행의 점과 두 막대, 맨 아래 행이 띠 안에 있는 것}

\slideheading{18}{대상별 전이 오차 변화}{0.6}

같은 라벨 0개 조건의 대상별 지도입니다. 원 하나가 대상 하나이고 파랑이 오차 감소입니다. 왼쪽 직접 ML은 알래스카 삽도를 비롯해 어두운 원이 많고, 오른쪽 증강은 대부분 흰색에 가깝습니다. 2~cm 넘게 나빠진 대상이 18개 대 2개입니다. 예외는 왼쪽 아래 티베트 고원입니다. 원천 계수 Stefan 오차가 약 245~cm인 심부 레짐이라 직접 ML이 60~cm 가까이 낫습니다. 계수가 크게 틀린 곳에서는 물리식이 출발점이 될 수 없으므로 티베트는 따로 보고합니다.

\pointat{알래스카 삽도의 어두운 원, 티베트 삽도의 짙은 파란 원, 컬러바 방향}

\begin{optbox}
두 신뢰구간이 모두 0을 넘은 대상으로 세면 14개 대 8개입니다. 알래스카 대상에서는 증강도 원천 계수 Stefan보다 1.79~cm 나쁘지만, 직접 ML의 35.3~cm를 16.3~cm로 낮췄습니다. 손해를 줄이는 것이지 없애는 것은 아닙니다.
\end{optbox}

\slideheading{19}{기존 ALT 지도와 물리 기준선}{0.5}

라벨 0개의 비교 기준을 기존 지도까지 넓혔습니다. 원천 계수 Stefan 대비 라벨 연도의 융해 도일을 쓴 연도 정합 Stefan은 1.00~cm 작았고, 기존 지도 CCI v5는 학습 지점 주변 5~km를 가린 뒤에도 23.8~cm 컸습니다. 맨 아래처럼 직접 ML은 연도 정합 Stefan보다 3.24~cm 컸습니다. 따라서 라벨이 없을 때의 출발점은 물리식입니다. 기존 지도 비교는 결과를 본 뒤 설계했고 원인은 단정하지 않습니다.

\pointat{끊어진 축의 CCI v5 행, 맨 아래 직접 ML 행}

\begin{optbox}
Stefan·위성 제품 앙상블은 두 가중의 판정이 달라 미결정입니다. Wei 2026 지도는 2지역에서 미결정이었고, 대상 지역 안 CALM 지점도 학습에 썼으므로 우열을 말하지 않습니다.
\end{optbox}

\slideheading{20}{라벨 10개 계수 재보정}{0.7}

세로축은 원천 계수 Stefan 대비 오차 변화입니다. 왼쪽 주 4지역 평균에서 파란 재보정 Stefan은 10개에서 2.45~cm 줄였습니다. 두 신뢰구간 모두 0 아래, 보정 뒤 $p$ 0.001의 사전 등록 판정입니다. 그런데 오른쪽 레나델타·캐나다 평균에서는 재보정 Stefan이 0선 위로 올라갑니다. 유의하게 준 지역은 러시아 서부 하나로 11.96~cm이고 캐나다는 2.26~cm 나빠졌습니다. 캐나다는 지역 평균 계수가 원천과 거의 같지만 라벨 위치마다 계수가 달라, 뽑힌 라벨로 맞춘 계수가 채점 블록에 맞지 않았다고 해석합니다. 사후 해석입니다.

\pointat{왼쪽 파란 선의 라벨 10개 점(−2.45), 오른쪽 파란 선이 0선 위로 올라가는 것}

\begin{optbox}
오른쪽의 직접 ML은 라벨 40개부터 0선 아래로 내려옵니다. 라벨이 수십 개를 넘으면 직접 ML도 쓸 자료가 생기므로 뒤에서 방법을 라벨로 고르는 규칙을 둡니다.
\end{optbox}

\slideheading{21}{물리 잔차 결합의 추가 이득}{0.5}

주황색이 재보정 Stefan 위에 잔차 가중 0.25의 잔차 학습을 얹은 물리 잔차 결합입니다. 라벨 10개의 재보정 대비 추가 이득은 0.18~cm였고, 보정 전 $p$ 0.034, 보정 뒤 0.136이라 차이를 확인하지 못했습니다. 전량에서는 0.40~cm 작았지만 0.5~cm 미만입니다. 따라서 라벨 10개까지 이득의 대부분은 재보정입니다. 사전 등록 판정입니다.

\pointat{왼쪽의 파란 선과 주황 선이 거의 겹치는 것, 오른쪽의 두 띠가 겹치는 것}

\begin{optbox}
표형 파운데이션 모델을 잔차 학습기로 쓰면 재보정 대비 라벨 3개에서 0.32, 10개에서 0.55~cm의 작은 추가 이득이 있었습니다. 잔차 구조의 이득은 학습기 종류에 묶이지 않습니다.
\end{optbox}

\slideheading{22}{대상 라벨 교차검증 선정}{0.6}

라벨이 수십에서 수백 개일 때입니다. 라벨 안 블록 5겹 교차검증으로 후보 7개 가운데 오차가 가장 작은 방법을 고르는 선정 규칙을 고정 잔차 레시피와 비교했습니다. 레나델타·캐나다 평균에서는 40개에서 0.87, 160개에서 1.32~cm 작았고 두 행 모두 0.5~cm 한계 안에서 비열등입니다. 다만 이득은 캐나다에서 왔습니다. 네모인 캐나다는 2.10과 2.72~cm 줄었고 마름모인 레나델타는 +0.35와 +0.07~cm입니다. 사후 분석입니다.

\pointat{아래 두 행의 점과 막대, 캐나다 네모와 레나델타 마름모}

\begin{optbox}
주 4지역의 10개와 전량은 미결정이고 10개에서는 비열등이 성립하지 않습니다. 40개의 0.87~cm는 보조 신뢰구간에서 판정이 달라 분할 독립 가정에 의존합니다. 후보 7개는 Stefan 3종, 잔차 가중 2종, 증강 결합, 물리 유사라벨 증강이며 10개 미만이면 재보정 Stefan을 씁니다.
\end{optbox}

\slideheading{23}{워크플로 순차 적용 결과}{0.6}

배치, 라벨 10개 진단, 방법 규칙을 순서대로 적용한 결과입니다. 얕은 레짐 5지역을 새로 나눈 분할에서 평가했습니다. 위 묶음은 무작위 배치에 고정 레시피를 쓴 경우 대비로 40개 미결정, 160개 0.93~cm 감소입니다. 아래 묶음은 재보정 Stefan 대비로 10개 동등, 40개와 160개에서 1.35와 1.51~cm 감소입니다. 160개 이득은 캐나다에서 왔고, 러시아 동부, 알래스카, 레나델타는 일부 라벨 수에서 0.5~cm 넘게 나빴습니다. 배치 단계는 무작위로 귀결됐고, 같은 지역을 다시 나눈 것이라 독립 검증이 아닙니다.

\pointat{위 묶음 두 행과 아래 묶음 세 행, 오른쪽의 옅은 지역 기호}

\slideheading{24}{지역 내 라벨 수 곡선}{0.5}

라벨이 많은 지역 안의 결과입니다. 블록 절반을 채점에 남기고 분할을 25회 반복했습니다. 세로축은 재보정 Stefan 대비입니다. 알래스카의 주황색 물리 잔차 결합은 라벨 1000개에서 0.54~cm 작았고 두 신뢰구간 모두 0 아래입니다. 보라색 직접 ML은 1000개에서도 0선 위입니다. 캐나다와 레나델타는 차이를 확인하지 못했습니다. 점선은 오차 하한으로 알래스카 11.31~cm입니다. 첫 시험이 기각된 뒤 설계를 바꾼 시험이라 사후 분석입니다.

\pointat{알래스카 패널의 라벨 1000개 주황 점, ‘Error floor’ 점선}

\begin{optbox}
RMSE로는 14.41에서 13.87~cm이고 오차 하한은 공변량이 같은 셀 사이의 실측 차이로 추정했습니다. RMSE 개선률 3.7\%는 하한 위의 줄일 수 있는 오차 제곱으로는 19\% 감소이고 대회 설정의 39\%와 같은 방향입니다. 부록 10에 있습니다.
\end{optbox}

\slideheading{25}{보정량의 공간 분포}{0.6}

알래스카에서 기계학습이 고친 오차를 공간 규모로 나눴습니다. 오른쪽은 오차를 ERA5 기후 격자 사이 성분과 격자 안 성분으로 나눈 것입니다. Stefan 예측은 격자 안에서 같은 값이므로 나눌 수 있습니다. 총 0.52~cm 감소 가운데 격자 사이는 1.03~cm 감소, 격자 안은 0.01~cm로 동등입니다. 오차 제곱 감소의 약 99\%가 기후 격자 단위 편향 보정입니다. Stefan 오차 제곱의 72\%가 격자 안 변동인데 기계학습이 설명한 몫은 1\% 미만입니다. 즉 격자 안 상세도는 더하지 못했습니다. 알래스카 지역 안에 한정한 사후 분석입니다.

\pointat{왼쪽 지도(블록별 오차 변화, 파랑이 감소), 오른쪽 ‘Between’ 열과 ‘Within’ 열의 알래스카 점}

\begin{optbox}
세 지역 모두에서 격자 안 설명 비율은 8\% 미만이고, 세 지역 평균에서 격자 안 예측은 오차를 0.09~cm 오히려 늘렸습니다. 격자 묶음을 토양 융해 도일로 바꾼 민감도에서도 알래스카 결과는 같았습니다.
\end{optbox}

\slideheading{26}{라벨 위치 전략}{0.5}

다음 라벨을 어디에 둘지입니다. 캐나다에서 라벨 100개를 셀 무작위로 뽑은 경우와 공변량 공간에서 덜 덮인 곳을 차례로 고른 경우입니다. 무작위는 서쪽 큰 블록에 몰리고 공변량 분산은 북쪽 섬까지 퍼집니다. 캐나다 지역 안 50에서 100개에서 분산 배치는 무작위보다 2.54에서 3.21~cm, 새 지역 적용의 40개에서 2.39~cm 작았습니다. 그러나 레나델타에서는 블록 층화가 100개에서 약 1.7~cm 키웠습니다. 효과가 지역에 따라 다르므로 사전 지정 절차는 무작위로 확정했습니다.

\pointat{왼쪽의 큰 회색 원 하나, 오른쪽의 북쪽 섬 주황 점들}

\begin{optbox}
레나델타는 라벨이 몇 블록에 몰려 있어 블록마다 고르게 뽑으면 계수 추정이 치우친다고 보지만 원인 분석은 하지 않았습니다. 예측 분산 기반 능동 선택은 알래스카와 레나델타에서 무작위보다 0.12에서 2.79~cm 나빴습니다.
\end{optbox}

\slideheading{27}{독립 지역}{0.5}

학습에 쓰지 않은 지역에서도 결론이 유지되는지입니다. 러시아 중부를 더한 얕은 레짐 5지역에서, 라벨 40개까지 직접 ML이 원천 계수 Stefan보다 유의하게 나은 지역은 0개였습니다. 오른쪽 위 두 행은 전량에서 물리 잔차 결합이 원천 계수 Stefan보다 5지역 평균 3.19~cm 작았다는 것이고, 아래 두 행은 라벨 10개의 재보정 대비 추가 이득으로 동등입니다. 티베트는 심부 레짐이라 따로 보고합니다. 확인적 독립 지역이 4에서 5개라는 것이 한계입니다.

\pointat{지도의 흰 원, 오른쪽 ‘러시아 중부’ 행의 긴 블록 등가중 막대(미결정)}

\slideheading{28}{예측 구간}{0.4}

보조 결과인 예측 구간입니다. 점선이 목표 0.90입니다. 알래스카에서 보정한 보통의 기계학습 구간을 새 지역에 그대로 쓰면 포함률이 0.32에서 0.73으로 떨어지고, 지역을 집단으로 보정한 계층 conformal 구간은 주 4지역 평균 0.86입니다. 보정 집단이 4개라 유한표본 보장은 없으므로, 새 지역에는 지역 단위 보정이 필요하다는 것까지가 결론입니다.

\pointat{‘알래스카 보정 구간’ 행의 긴 범위 선, ‘지역 단위 보정’ 행의 큰 점}

\slideheading{29}{알래스카 1 km ALT 지도}{0.7}

지도입니다. 라벨이 많은 지역의 최종 방법인 물리 잔차 결합으로 알래스카 영구동토 구역 표시 셀 858,517개를 그린 1~km 지도이고 250~m 음영 지형 위에 얹었습니다. 정확도는 0.5° 블록 5겹 교차검증 값만 말씀드립니다. 물리 잔차 결합 13.80~cm, 재보정 Stefan 14.30~cm이고 90\% 예측 구간은 상수 폭 ±20.2~cm입니다. 세부 무늬는 지형 입력에서 오지만 정보 해상도는 기후 입력의 약 10~km에 묶입니다. 지도 셀의 73\%는 공변량이 학습 범위 밖이라 외삽으로 읽어야 합니다.

\pointat{북쪽 해안의 옅은 띠와 내륙의 짙은 띠, 컬러바 40–70 cm}

\begin{optbox}
외삽 셀에서는 구간의 포함률이 보장되지 않습니다. 범위 밖 공변량은 주로 지형 거칠기, 경사, TPI이며 라벨 343곳의 지형 범위가 전체 영역을 덮지 못하기 때문입니다. 세부 패널은 부록 12에 있습니다.
\end{optbox}

\slideheading{30}{라벨 수에 따른 지도 변화}{0.7}

같은 지역의 1~km 지도가 라벨 수에 따라 어떻게 바뀌는지를 레나델타로 보입니다. a는 라벨 0개의 원천 계수 Stefan입니다. b는 라벨 10개의 재보정 Stefan인데 계수가 1.62에서 2.04로 커져 전체가 깊어집니다. c와 d는 40개와 160개의 물리 잔차 결합으로, 잔차 가중 1.0이 선택되어 북동부가 뽑힌 라벨에 따라 크게 달라집니다. e의 전량 3,037개에서 안정됩니다. 요지는 라벨이 적을 때 지도는 어느 라벨이 뽑혔는지에 좌우되고 수백 개부터 안정된다는 것입니다. 정확도는 말하지 않으며, 라벨 40개 지도가 정확하다는 뜻이 아닙니다.

\pointat{윗줄 ALT, 아랫줄 전량 대비 변화. b 의 짙어진 색, c·d 의 북동부 짙은 영역, j 의 거의 흰 지도}

\slideheading{31}{표시 해상도와 정보 해상도}{0.6}

1~km 지도가 1~km만큼의 정보를 담는지 확인했습니다. 같은 지도를 왼쪽은 1~km 셀 값, 오른쪽은 0.1° 셀 평균으로 그렸습니다. 0.1° 평균이 1~km 지도 분산의 알래스카 97\%, 레나델타 61\%를 설명했습니다. 그래서 1~km는 표시 해상도이고 정보 해상도는 기후 입력 ERA5-Land 0.1°, 약 10~km와 토양 입력 약 5~km에 묶입니다. 앞의 이득 분해와 같은 방향이며, 그래서 고해상 ALT 지도라고 말하지 않습니다.

\pointat{a 와 b 의 같은 띠 구조, c 와 d 의 차이}

\begin{optbox}
0.1° 셀 안 1~km 값의 표준편차는 알래스카 1.8~cm(전체 9.5), 레나델타 1.1~cm(전체 1.8)입니다. 알래스카의 ERA5-Land 0.1° 격자는 약 4–6 × 11~km입니다.
\end{optbox}

\slideheading{32}{기존 ALT 제품과의 차이}{0.5}

우리 지도와 기존 제품이 어디서 얼마나 다른지입니다. 정확도 비교가 아니며 정확도는 같은 채점 셀 비교로만 판단했습니다. 영역 평균은 물리 잔차 결합 56.1~cm이고 CCI v5 105.5, Wei 2026 70.8, Aalto 2018 78.8, Yi·Kimball 105.1~cm로 기존 제품이 모두 깊습니다. 공간 상관은 재보정 Stefan과 0.96, 기존 제품과는 −0.15에서 0.46입니다. 차이가 큰 곳은 알래스카 내륙 등이고 원인은 단정하지 않습니다.

\pointat{a 의 mean 56.1 과 b–e 의 mean·r 값, g 의 올리브색 내륙}

\slideheading{33}{레나델타 입체 보기}{0.5}

레나델타 1~km 지도를 250~m 지형 위에 얹어 남동쪽에서 본 모습입니다. 지형은 열두 배 과장했고 남서쪽 구릉의 최고 높이는 약 580~m, 델타 평원은 약 30~m 아래이며 색 범위는 35에서 50~cm입니다. 평원의 ALT 무늬는 지형의 수로나 호수를 따르지 않고 기후 격자 입력의 띠를 따릅니다. 앞 쪽의 정보 해상도 결론을 눈으로 확인하는 그림이고, 축이 없으므로 모양과 상대 기복만 읽습니다.

\pointat{왼쪽 아래 구릉의 기복, 평원을 가로지르는 띠 무늬}

\slideheading{34}{결과 종합}{0.9}

결과를 라벨 수 구간별로 종합합니다. 왼쪽 주 4지역 평균에서는 재보정과 저가중 잔차 결합이 함께 내려가 10개에서 약 2.5~cm 줄었고 직접 ML은 0선 위입니다. 오른쪽 레나델타·캐나다 평균에서는 캐나다의 재보정 오차 때문에 라벨을 쓰는 방법들이 0선 위에 놓이지만, 짙은 갈색 점인 교차검증 선정은 40개부터 그 아래로 내려옵니다. 라벨이 없으면 물리식이며 기존 지도 제품보다 23.8~cm 정확했습니다. 10개면 재보정만으로 2.45~cm를 줄였고 잔차 학습의 추가 이득은 확인되지 않았습니다. 40개 이상에서는 교차검증 선정과 잔차 결합이 재보정 Stefan보다 작았고 160개에서 1.5~cm였습니다. 이것이 새 지역에서 학습, 평가, 추가 관측의 순서를 정하는 워크플로의 근거입니다.

\pointat{a 의 라벨 10개 점, b 의 짙은 갈색 점 세 개, 아래 권고 띠}

\begin{optbox}
160개의 1.5~cm는 같은 지역을 다시 나눈 풀 평균입니다. b 풀에서 라벨을 쓰는 방법이 0선 위에 놓이는 것은 캐나다의 재보정이 원천 계수보다 나빴기 때문이고, 그래서 40개부터는 고정 레시피가 아니라 대상 라벨로 고르는 규칙을 권고합니다.
\end{optbox}

\slideheading{35}{연구 의의와 차별성}{0.6}

의의는 셋입니다. 첫째, 지역 홀드아웃과 100~km 완충, 라벨 0개부터 전량까지의 곡선, 같은 라벨로 재보정한 물리 기준선을 함께 갖춘 ALT 연구는 조사 범위에서 찾지 못했습니다. 방법 하나하나는 새롭지 않고 기여는 평가 설계와 음성 결과를 포함한 결과, 워크플로입니다. 둘째, 무작위 분할은 직접 ML의 새 지역 오차를 절반으로 과소평가합니다. 셋째, 규모입니다. 평가 지역 7곳, 전이 대상 27개, 라벨 수 8단계, 방법 12종과 학습기 10종, 실행 전 등록한 가설 174개, 모형 적합 최소 863,229건입니다.

\pointat{왼쪽 표의 ‘등록 가설 174개’와 ‘모형 적합’, 오른쪽 격자의 빈 칸}

\begin{optbox}
초록에 판정어를 쓸 수 있는 사전 등록 대비 10개 가운데 다중 비교 보정 뒤 방향이 확정된 것은 6개, 동등 1개, 보정 전만 유의 1개, 미결정 2개입니다.
\end{optbox}

\slideheading{36}{주장과 근거}{0.8}

주장 다섯 가지입니다. 첫째, 라벨 없는 새 지역에서는 물리식이 출발점입니다. 근거는 물리 유사라벨과 섞은 라벨의 차이 1.6~cm와 기존 지도 CCI v5의 23.8~cm입니다. 둘째, 라벨 몇 개의 이득은 계수 재보정입니다. 재보정 2.5~cm, 추가 ML 이득 확인되지 않음, 사전 등록입니다. 셋째, 수십 개부터는 교차검증 선정이 고정 레시피 이상입니다. 0.9에서 1.3~cm, 캐나다 주도, 사후 분석입니다. 넷째, 수백 개에서 잔차 ML의 이득은 약 0.5~cm이고 그 99\%는 기후 격자 사이 성분입니다. 다섯째, 원천 확대, 배치 규칙, 학습기 교체, 고해상 입력은 이득을 더하지 않았습니다. 그래서 병목은 모델이 아니라 정보입니다.

\pointat{표의 다섯 행, 오른쪽 ‘근거 종류’ 열}

\begin{optbox}
표의 ‘안전’은 비열등 시험 통과가 아니라 큰 실패를 막는 출발점이라는 뜻으로 한정합니다. 다섯째의 근거는 물리식과 제품의 적층, 알래스카에서 고른 배치 규칙의 이전, 학습기 10종, 10에서 500~m 입력 시험이 모두 이득을 더하지 않았다는 결과입니다.
\end{optbox}

\slideheading{37}{한계와 해석}{0.7}

한계 다섯과 대응입니다. 첫째, 학습 원천의 78에서 94\%가 알래스카입니다. 지역 홀드아웃으로 외삽을 평가했지만 확인적 독립 지역이 4에서 5개라는 사실은 남습니다. 둘째, 배치 알고리즘이 무작위로 귀결됐습니다. 분산 배치 이득이 캐나다 −2.5, 레나델타 +1.7~cm로 지역에 따라 다르다는 사실 자체가 지침입니다. 셋째, 격자 안 이득이 0.15~cm로 작습니다. 정보 해상도의 결론입니다. 넷째, 워크플로 검증이 같은 지역 재분할이라 독립 검증이 아니며 새 지역 결과는 10월 11일 뒤 확정됩니다. 다섯째, 지역 안 시험은 사후 분석이므로 사전 등록 시험과 분리해 표시합니다.

\pointat{표의 왼쪽 ‘한계’ 열 다섯 항목}

\slideheading{38}{기대효과와 향후 계획}{0.5}

기대효과는 라벨 투자 판단 근거, 방법 선택 지침, 평가 표준, 지도의 정보 해상도 명시의 넷입니다. 계획은 셋입니다. 격자 안 입력 시험은 WorldCover 10~m 피복 비율과 Sentinel-2 20~m 식생 지수로 먼저 했는데, 격자 안 오차는 0.15~cm 줄어 0.5~cm 기준에 못 미쳤고 레나델타 전체 오차는 1.5~cm 줄었습니다. 위성 적설 자료는 10월 8일까지 들어옵니다. 독립 지역은 극지연구소와 KPDC 자료 협업으로 확보하고, 논문은 Scientific Reports에 제출합니다.

\pointat{오른쪽 표의 ‘격자 안 입력 시험’ 행}

\slideheading{39}{Thank you}{0.2}

이상으로 마치겠습니다. 질문 주시면 답변드리겠습니다. 부록에 학습기별 결과, 판정 단어의 정의, 지도 세부, 추가 실험 요약이 있습니다.


% ---------- 2부 ----------
\partpage{2부}{발표자 배경 지식}{말하지 않는 내용 · 질문에 답할 때 참고 · 문어체}
\markboth{2부 발표자 배경 지식}{}
% part2_background.tex : 발표자 배경 지식(말하지 않는 내용). presentation_script.tex 에서 \input 한다.
% 출처: docs/EXPERIMENT_PLAN_FINAL_BATCH_2026-10-04.md(0.3, 1절, 8절), 추가 등록 XK/XL·XM, docs/MORNING_REPORT_2026-10-05.md,
% paper/manuscript/en/sections/methods.tex, paper/claims/*/README.md, paper/registry/experiments.csv, docs/RESEARCH_OVERVIEW_2026-10-02.md,
% docs/QA_FINAL_REVIEW_2026-10-02.md, docs/QA_FOLLOWUP_2026-10-04.md, 코드 모듈 머리 설명(scripts/, src/polar/, tools/, configs/rescale/).

\section{연구 수행 경과}

\subsection{시간순 경과}

\begin{center}\tabsmall
\begin{tabularx}{\textwidth}{@{}>{\raggedright\arraybackslash}p{2.4cm}Y@{}}
\toprule
시기 & 한 일과 결과 \\
\midrule
2026-06-30 – 07-13 & 2026 극지 빅데이터·인공지능 활용 경진대회(주최 극지연구소, 주관 KPDC) 준비. ERA5-Land 재채점으로 지역 전이 오차 108.5 → 87.3 cm(WorldClim 대비 20\% 감소). 공간 딥러닝 6모델 비교에서 앙상블 16.95, GBM 17.24 cm 동률. ‘17 cm 물리 하한’, ‘12.97 cm SOTA’ 문장은 07-06 에 폐기(범위 축소 아티팩트). \\
07-14 – 07-21 & 지역 홀드아웃(LORO)에서 Stefan 18.2 cm, 순수 ML 40.6 cm, 잔차 GBM 48.8 cm. 다지역 통합 학습은 결측 처리 아티팩트로 미채택. ‘물리 = 라벨 공급원(백본)’으로 역할 정리. \\
07-23 – 07-30 & 대회 실험 S0–S12: 알래스카 지역 안 Stefan 14.46, 물리 잔차 결합 13.33 cm(같은 교차검증에서 고른 최선 설정), CatBoost 15.61 cm. 전이에서는 Stefan 앵커 21.26 cm 가 최선. 07-30 예선 분석 보고서 제출(KPDC 콘슬 지중온도·코어 자료를 알래스카 지역 안 점검증으로만 사용). \\
08-26 – 08-31 & 본선 덱 v5 확정, 논문화 판정(중위권 게재권), 투고지 Scientific Reports 로 방향(08-31 기록은 CRST 1순위와 Sci Rep 확정이 엇갈림. 09-08 이후 모든 계획은 Sci Rep 기준). \\
09-08 – 09-21 & 사전 등록 실험 E1–E4: 대회 헤드라인 24.11 → 22.92 cm 는 중첩 재평가에서 25.49 cm 로 나와 철회(승자의 저주). 기존 주장 8건 감사(뒤집힘 0, 하향 3). M1 마스터 요인 실험 등록·실행, 재현 게이트(0.1 cm) 통과. \\
09-22 – 09-26 & 전이 수정 H18–H24(정보·계수·학습 목표 축 모두 기각), 라벨 예산 H25–H30, 09-26 감사(지적 22건 가운데 21건 확정: 행 단위 iid CI 폐기 → 블록 부트스트랩, ‘ρ −0.84’ 폐기 등), 최종 계획 F1–F11 실행(지지 3, 부분 3, 기각 5). \\
09-29 – 09-30 & 연구 틀 재정의: 대상 라벨 수 $n$에 대한 오차 변화 곡선, 두 물리 기준선(원천 계수, 같은 라벨 재보정). 신규성 입장 정리(기법 신규성 없음, 평가 설계·음성 결과·워크플로가 기여). LG(라벨 수 격자) 사전 등록(커밋 40be64c 15:45:46) 뒤 Rescale 실행. LGX·LGT·LGD·LGU·LGF·WRAPUP 등록, 09-30 07:20 첫 결과 회수, 판정 기록 J1–J7. \\
10-01 – 10-02 & 사용자 지정 주제: 물리 증강·잔차 ML 의 라벨 수별 능력과 새 지역 워크플로. 지역 안·워크플로 시험 WF0–WF10(결과 열람 뒤 설계, 실행 전 등록). 10-02 연구 개요와 원고 재구성안, WF9 토양 격자 민감도. \\
10-04 & 추가 실험 XA–XJ 사전 등록(T0 = 커밋 2678100, 14:22:39). XA 봉인 표 첫 열람 21:07. \\
10-05 & Rescale 작업 A(elm, 02:18–02:38, 약 1.40 달러)와 로컬 실행으로 XB·XC·XD·XE 1단계·XG·XH·XI·XJ 열람(03:18–05:43). XK·XL 등록·결과(지도·격자 크기). XM 등록(13:23)·열람(14:20). 제목 확정, 덱 2차 개정, 원고 영문 39쪽·SI 76쪽·국문 29쪽 조립. \\
\bottomrule
\end{tabularx}
\end{center}

\subsection{틀의 변화}

대회 제출 제목은 ‘희소 관측 조건에서 물리경험식 유사라벨 증강과 기계학습을 이용한 영구동토 활동층 두께 예측 및 독립 관측 결합에 의한 지역 간 전이’였고, 주장은 (1) 유사라벨 증강의 소수 라벨 지역 개선(캐나다 34.0 → 31.8 cm), (2) 전이에서 물리식 단독을 넘는 구조 없음과 위성 제품 결합의 편향 상쇄(24.11 → 22.92 cm), (3) 보정 예측구간 포함률 44.6 → 93.4\% 였다. 이 가운데 (2)의 22.92 cm 는 09-08 중첩 재평가에서 철회되었고, 증강 이득은 라벨 0개 조건으로 한정되었다.

09-29 에 평가 축을 ‘대상 라벨 수 $n$에 대한 오차 변화 곡선’ 하나로 통일하고 기준선을 원천 계수 Stefan 과 같은 라벨로 재보정한 Stefan 둘로 정했다. 09-30 판정 직후의 주제(‘물리식이 강하다, 무작위 분할은 과대평가한다’)는 사용자 검토에서 ‘너무 뻔하다’는 지적을 받았고, 10-01 에 주제를 ‘물리 증강·잔차 ML 의 라벨 수별 능력과 새 지역 워크플로(학습·평가·추가 탐사)’로 다시 정했다. 무작위 분할 과대평가는 보조 근거로 내렸다. 현재 제목은 10-05 사용자 확정 문구다.

\subsection{사전 등록 실무}

\begin{itemize}[leftmargin=1.4em, itemsep=0.15em]
  \item \textbf{적합 전 커밋}: 가설, 대비, 판정 규칙, 허용 문장, 본문·SI 배치를 git 에 커밋한 뒤 실행한다. 커밋 시각이 T0 이고 마감은 T0 기준으로 센다(예: XE 2단계 T0 + 96 h, XF 새 지역 T0 + 7일). 제3자 타임스탬프는 없고 커밋 시각은 KST 다.
  \item \textbf{봉인 산출}: 집계기는 판정 표를 \pth{data/processed/xbatch/<이름>/sealed/} 에 쓰고 화면에는 파일 이름, 행 수, sha256 앞 16자만 남긴다. 표준 출력은 \pth{grep -v -E '판정|verdict|Δ|delta|rmse|...'} 로 거른다. 열기 전에 \pth{sealed\_manifest.json} 의 sha256 과 대조하고, 첫 열람 시각을 계획서 8절에 적는다.
  \item \textbf{열람 순서}: 알래스카 계열 배치 표 → 작업 R2a 제출 뒤 나머지 판정 표 → 배치 정책 동결 커밋 뒤 비알래스카 행 → 시험 표 열람 뒤 교차검증 자료. 재현 관문을 통과하기 전에는 판정 표를 열지 않는다.
  \item \textbf{재현 관문}: 블록별 SSE 를 기준 조각과 대조한다. 같은 노드 종류 |ΔSSE| 0, elm–hematite 사이 0(11,814키), 로컬–Rescale 사이 ≤ 1e-4 cm² 또는 상대 1e-9(기록된 최대 6.8e-6 cm²). 관문을 넘으면 판정을 멈추고 원인을 개정 이력에 적는다. 10-05 XE 는 처음에 잘못된 수준으로 채점되어 실패로 표시되었다가 등록 수준으로 다시 채점해 35,442키 모두 통과했다.
  \item \textbf{판정 규칙}: Δ = RMSE(A) − RMSE(B). 채점 블록 재표집 10,000회(LG 본 판정 L1–L8 과 예측 구간 실험은 1,000회)로 셀 가중·블록 등가중 두 95\% CI. 우세(오차 감소) = 두 CI 상한 < 0, 열세 = 두 하한 > 0, 동등 = 네 끝값 절댓값 ≤ 0.5 cm, 미결정 = 그 밖. 보조 δ 1.0 cm 와 0.02 × 원천 Stefan RMSE 를 병기하고 판정이 다르면 ‘한계 의존’. 분할 독립 가정 CI 와 다르면 ‘분할 독립 가정 의존’. 비열등 = 두 상한 < +0.5 cm. Holm 보정은 초록 대비 10개(m 10), XC(m 8), XG(m 6)에서 문구를 정하고 나머지는 보조 열. 점 추정 절댓값 0.5 cm 미만의 우세·열세는 ‘통계적으로 구별되나 크기는 0.5 cm 미만’.
  \item \textbf{표지}: 설계 표지(사전 등록, 결과 열람 뒤 설계, 사후 분석, 사후 설계)와 맹검 표지(맹검, 비맹검 부분 포함, 재현(비맹검), 등록 시점에 결과 존재(미열람))를 가설마다 적는다. 해석 문장은 다섯 갈래(우세·열세·동등·미결정·판정 불가)를 결과 전에 고정한다.
\end{itemize}

\section{코드 구성}

\subsection{언어와 환경}

\begin{center}\tabsmall
\begin{tabularx}{\textwidth}{@{}>{\raggedright\arraybackslash}p{3.0cm}YY@{}}
\toprule
항목 & 로컬(공유 서버) & Rescale(추가 묶음 고정판) \\
\midrule
Python & 3.9.18 & 3.11.10 \\
CatBoost & 1.2.10 & 1.2.10 \\
scikit-learn / pandas / numpy / scipy & 1.3.0 / 2.1.4 / 1.26.4 / 1.11.4(\pth{environment-lock.txt}) & 1.9.1 / 2.3.3 / 2.1.1 / 1.17.1(\pth{scripts/rescale/xbatch\_constraints.txt}) \\
torch & 2.6.0+cu124(DeepSets 배치 정책 학습에만) & 2.4.1(deps.log) \\
Cartopy & 0.23.0(지도) & 쓰지 않음 \\
조판 & xelatex, Pretendard, kotex & \\
하드웨어 & 128코어 공유, 우리 상한 32스레드·nice 10, GPU 0–4 는 DeepSets 전용 & elm 96코어(AMD EPYC Milan, 3.24 달러/시간), hematite 64코어(Milan-X, 4.22), iolite-4 64코어 + T4 4기(5.70, LG·LGX) \\
\bottomrule
\end{tabularx}
\end{center}

\subsection{폴더 구성}

{\tabsmall
\begin{longtable}{@{}>{\raggedright\arraybackslash}p{4.3cm}>{\raggedright\arraybackslash}p{11.3cm}@{}}
\toprule
폴더 & 역할과 대표 파일 \\
\midrule
\endfirsthead
\toprule
폴더 & 역할과 대표 파일 \\
\midrule
\endhead
\bottomrule
\endfoot
\pth{scripts/0\_download/}(10) & 원자료 내려받기: Copernicus DEM 타일, ERA5-Land 월평균 2010–2024, SoilGrids 창, ReSALT InSAR, XE 2단계 취득. \\
\pth{scripts/1\_data\_prep/}(62) & 라벨 표와 공변량: \pth{build\_fidelity\_base\_v3.py}(주 표), \pth{v4}(새 지역), \pth{parse\_ext\_*.py} 22개(외부 라벨 원천별 파서), 지도 격자 \pth{build\_map\_grid\_\{lena,alaska\}\_v1.py}, XM 특징 \pth{xm\_landcover\_s2\_features.py}. \\
\pth{scripts/2\_evaluation/}(61) & 분석·채점·지도 예측: 동결 모듈 \pth{h41\_validation\_ladder.py}(검증 사다리), \pth{xa\_c2\_gain\_decomposition.py}, \pth{xk\_support\_scale\_v1.py}, \pth{xl\_map\_products\_v1.py}, \pth{map\_alaska\_v1.py}, \pth{h49\_transfer\_map.py}(레나), \pth{label\_sequence\_v1.py}, \pth{wf0\_reanalysis.py}. \\
\pth{scripts/3\_deep\_learning/}(88) & 실험 하네스: 동결 모듈 \pth{h40\_label\_grid.py}(라벨 수 격자), \pth{h42\_label\_grid\_ext.py}(확장·판정·풀), \pth{h54\_workflow.py}(지역 안·워크플로), 공용 골격 \pth{xbatch\_core.py}, 추가 묶음 \pth{x\_*.py} 11개. 이름은 역사적이며 최종 묶음에 신경망 학습기는 없다. \\
\pth{scripts/4\_visualization/paper\_v3/}(19) & 논문 그림 \pth{fig1.py}–\pth{fig7.py}, \pth{table1.py}, \pth{si\_figs.py}, 지도 \pth{maps\_alt\_v3.py}, \pth{hillshade\_maps\_v1.py}, \pth{lena\_3d\_v1.py}, \pth{cover\_bg\_v1.py}. 덱은 \pth{deck/build\_paper\_report.py}. \\
\pth{scripts/rescale/}(16) & 묶음 생성 \pth{make\_payload*.sh}, 노드 실행 \pth{run\_\{lg,lgx,wf,xbatch\}.sh}, 판 고정 \pth{xbatch\_constraints.txt}. \\
\pth{src/polar/}(31) & 공용 패키지: \pth{m1\_core.py}(\pth{load\_base}, \pth{half\_split\_blocks}, \pth{eval\_mask}), \pth{h4\_common.py}(\pth{seed\_of}, \pth{boot\_delta\_blocks}, \pth{boot\_weights}), \pth{fidelity.py}(공변량 25종 정의), \pth{cv\_schemes.py}(kNNDM), \pth{physics.py}(Stefan 계열 식 5종). \\
\pth{tests/}(31, 함수 626) & 합성 자료 단위 시험, 누설 시험(b·d·m 형식), 동결 해시·고정판·출력 필터·봉인 쓰기 시험. \\
\pth{tools/}, \pth{configs/rescale/}(16) & \pth{rescale\_client.py}(REST: upload → job → submit → poll → fetch, 비용 상한·저우선순위 차단·장부 기록), 노드 종류별 YAML. \\
\pth{data/processed/xbatch/<이름>/} & \pth{shards/}(단위별 \pth{runs.csv}, \pth{blocksse.npz}, \pth{unit.json}), \pth{inputs/}, \pth{sealed/}(판정 표·매니페스트). \\
\pth{results/rescale*/} & 작업별 회수 산출(상태 CSV, 로그, 자료). \pth{results/rescale/\_ledger.csv} 가 지출 장부. \\
\pth{paper/} & \pth{registry/experiments.csv}(66행), \pth{claims/}(C1–C8, D, SI 2개: README + 표 사본 + sha256 매니페스트), \pth{manuscript/\{en,ko\}}. \\
\pth{outputs/figures/paper/v3\_restructure/} & Fig1–7, Table1 의 PDF·PNG·설명문·Source Data·값 점검 기록, \pth{maps/}, \pth{si/}. \\
\end{longtable}}

\subsection{동결 모듈과 하네스}

동결 모듈 네 개의 sha256 앞 16자는 다음과 같고, \pth{xbatch\_core.py} 가 불러올 때 대조해 다르면 멈춘다.

\begin{itemize}[leftmargin=1.4em, itemsep=0.05em]
  \item h40 \pth{7098f59dabe2b73f}, h42 \pth{22e215e435e157be}
  \item h54 \pth{cf9a1f6d0929c892}, h41 \pth{492373e4d37b5ea4}
\end{itemize}

h40 은 라벨 추출 \pth{draw\_cells}(seed = \pth{seed\_of(target, mode, split, n, draw)}), 대비 \pth{contrast}; h42 는 교차검증 겹 \pth{cv\_folds\_of}, 4분 판정 \pth{verdict4}, 보조 CI \pth{boot\_delta\_common}, 층화 평균 \pth{pool\_rows}, 지역 추론 \pth{region\_inference}; h54 는 최원점 순서 \pth{kcenter\_order} 와 지역 안·워크플로 단위; h41 은 검증 사다리(학습·채점 셀이 겹치지 않음을 단언)다. \pth{xbatch\_core.py} 는 20개 절로 된 공용 골격(실행 보호, 출력 제한, 분할 계획, 추출, 채점·관문, 대비·CI, 판정·Holm·문장 조립, 조각 쓰기, 봉인, 누설 시험, 묶음 매니페스트, 결정론 설정)이며 가설과 판정 규칙은 각 \pth{x\_*.py} 가 정한다.

\subsection{실험 실행 순서}

\begin{enumerate}[leftmargin=1.8em, itemsep=0.1em]
  \item \textbf{세기}(\pth{--count-only}): 분할 구조와 단위 수만 센다. 라벨 값을 쓰지 않는다.
  \item \textbf{스모크}(\pth{--smoke}): 평가에 쓰지 않는 분할에서 합성·소규모 실행. OMP 스레드 2, nice 10, taskset 코어 4개, grep 필터.
  \item \textbf{본 실행}: \pth{WF\_RESCALE=1} 또는 \pth{--allow-local} 이 없으면 거부. 워커 × 스레드(Rescale elm 22 × 4), \pth{--resume} 는 조각 세 파일과 설정 해시가 맞을 때만 건너뛴다.
  \item \textbf{재현 관문}(\pth{--gate}, \pth{--gate-only}, \pth{--gates-only}, \pth{--gate-wf9}): 기준 조각과 블록 SSE 대조.
  \item \textbf{집계·봉인}(\pth{--summarize-only}): 재표집 10,000회, 판정 표를 \pth{sealed/} 에 쓰고 매니페스트 기록.
  \item \textbf{열람}: 0.3 열람 순서대로 사람이 연다. 열람 시각과 해시를 8절에 적는다.
\end{enumerate}

시험은 \pth{CUDA\_VISIBLE\_DEVICES= OMP\_NUM\_THREADS=2 nice -n 10 taskset -c <4코어> python3 -m pytest -q tests/<파일>} 로 돌리고, 메모리는 \pth{prlimit --as}(또는 \pth{--data}) 10 GB 로 막는다. 재현성은 seed(\pth{seed\_of} = CRC32 mod $2^{31}$), CatBoost CPU 결정론(\pth{random\_seed}, \pth{thread\_count} 명시), 코드·설정·입력 표의 sha256, 묶음 정보의 git 커밋 해시로 확보한다. 물리식과 CatBoost 는 Rescale 과 로컬 사이 최대 차 3.6e-14 cm(34,360키), 노드 종류 사이 비트 단위 동일(11,814키)이었고, 신경망은 환경 간 최대 5.6 cm 달라(FT-Transformer 288키 가운데 43키가 0.5 cm 초과) 최종 묶음에서 제외했다.

\subsection{Rescale 사용과 비용}

장부(\pth{results/rescale/\_ledger.csv}, 13건)의 정산 기록은 LG·LGX 6건 220.26 달러(LG 본 실행 109.25, LGX A 63.37, B 45.97 등)다. WF 1–3차와 작업 A 등 7건(elm·hematite)은 장부에 실제 비용이 비어 있고 10-02 기록은 WF 5건을 약 2.98 달러로 합산해 223.24 달러를 적었다. 작업 A(10-05, elm 0.43 h)는 약 1.40 달러다. 누적은 약 225 달러이며 이 합계를 적은 문서는 없다. 누적 지출 + 제출 예정 최악 비용이 50 달러를 넘는 제출은 사용자 확인 뒤에만 한다. 10-05 의 나머지 실행(XC, XD 시험, XE 1단계, XG 셀 단위, XB 민감도, XM, 지도)은 로컬 CPU 와 GPU 0·1 에서 했다.

\section{자료}

\subsection{라벨}

{\tabsmall
\begin{longtable}{@{}>{\raggedright\arraybackslash}p{3.2cm}>{\raggedright\arraybackslash}p{12.4cm}@{}}
\toprule
항목 & 내용 \\
\midrule
\endfirsthead
\toprule
항목 & 내용 \\
\midrule
\endhead
\bottomrule
\endfoot
표의 판 & v1 17,423행(대회) → v2 17,572행(CALM 149행 추가, 중복 8 제거) → v3 17,572행(출처 표지 68행 변경, 본 실험) → v4 18,088행(새 지역 516행: 직접 329, 지온 유도 39, 미상 148). 주 자료 \pth{data/processed/fidelity\_base\_v3.csv}. \\
직접 라벨 셀 & 17,467개(탐침 또는 GPR). 셀 = 좌표 소수 넷째 자리 반올림(ABoVE, ALLena) 또는 CALM 지점, 값은 관측 평균. CALM 68셀(지온·융해관 메타)은 라벨로 쓰지 않음. 표 동결 뒤 찾은 지온 유도 셀 20개(알래스카 15, 캐나다 4, 티베트 1)는 표지만 달고 직접 라벨로 유지. \\
원천 & GTN-P/CALM(Streletskiy 2025), ALLena(Veremeeva 2025), ABoVE v2(Moore 2025). 러시아 중부·티베트는 공개 보관소 1 km 셀 평균. \\
지역별 행 / 1 km 위치 / 블록 & 레나델타 3,037 / 201 / 20, 캐나다 750 / 86 / 39, 러시아 서부 31 / 29 / 21, 러시아 동부 30 / 30 / 21, 알래스카 13,606 / 343 / 74, 러시아 중부(확충) 57 / 57 / 13, 티베트 132 / 132 / 34. 1 km 위치당 행 수는 알래스카 39.7, 레나 15.1, 캐나다 8.7. \\
연도 & 1990–2024. 예측은 2015–2020 기후값에 대한 정적 예측. 연별 변동은 공변량으로 예측되지 않아(R² ≤ 0.0015) 다년 평균 규약. \\
약관 판 & 재배포 약관이 확인되지 않은 원천(calm\_web\_subsites, cusp\_v1\_1, nsidc\_ggd353\_thawtube, ru\_yamal\_walker)을 뺀 ‘약관 확인분 판’이 주 판정. 이 때문에 그림·주 판정에서 뺀 셀은 캐나다 38, 북대서양 19, 러시아 서부 1. 북대서양(41행, 블록 13)은 전체 판에서만 적격이라 SI. 지온 유도 보조 행 원천 쪽(XJ-1·2)은 약관 근거 미기록으로 시험하지 않음. 계정이 필요한 보관소(TPDC 등)는 쓰지 않음. \\
KPDC & 대회 때 콘슬(수어드반도) 지중온도 프로파일 19종·2022 코어 18개·AWS 를 알래스카 지역 안 점검증으로만 사용(식별자 KOPRI-KPDC-00002125, 00002707, 00002955). 논문에서는 SI 사례·한계 서술. 10-04 극지연 김승희·정윤택 박사와 협업 예정. \\
\end{longtable}}

\subsection{공변량 25종과 해상도}

\begin{center}\tabsmall
\begin{tabularx}{\textwidth}{@{}>{\raggedright\arraybackslash}p{2.6cm}>{\raggedright\arraybackslash}p{3.6cm}Y@{}}
\toprule
군 & 원자료 & 열 \\
\midrule
지형 6 & Copernicus DEM GLO-30(30 m) & 고도, 경사, 향의 sin·cos(점 화소), TPI, 거칠기(33 × 33 화소 창, 약 1 km) \\
기후 8 & ERA5-Land 0.1°(약 10 km) 2015–2020 월평균 & 연평균 기온, TDD, FDD, √TDD, 최난월·최한월 기온, 1층 토양 온도, 적설 수당량 \\
토양 9 & SoilGrids 2.0 250 m 를 약 5 km 로 추출 & 점토·모래·실트, 용적밀도, 조립질, pH(5–15 cm), 유기 탄소 3층 \\
CCI 2 & ESA CCI 영구동토 ALT v4.0, 1 km & 1997–2021 평균, 유효 표지(같은 재분석으로 구동한 모형 출력이라 관측으로 보지 않음) \\
\bottomrule
\end{tabularx}
\end{center}

Stefan 입력 TDD = Σ max(월평균 기온, 0) × 일수, $s = \sqrt{\mathrm{TDD}}$(기온 기준). 채점 셀은 라벨, CCI 값, 1층 토양 온도 TDD 가 모두 유효한 셀이며 모든 방법이 같은 채점 셀을 쓴다. 라벨에서 유도한 양(관측 수, 연도)은 입력에서 제외하고 단위 시험이 이를 강제한다. 정보 해상도는 기후 약 0.1°, 토양 약 5 km 다.

\subsection{지도 격자와 추가 특징}

\begin{itemize}[leftmargin=1.4em, itemsep=0.15em]
  \item 알래스카: 영역 셀 899,633, 표시 셀 858,517(CCI 영구동토 비율 ≥ 50\%), 59–71.5°N·168–141°W, 1 km, EPSG:3338. 재보정 계수 1.6166(원천 1.5046, 원천 셀 3,860), 잔차 가중 0.5(블록 교차검증 선택), 블록 5겹 교차검증 RMSE 13.80(물리 잔차 결합) 대 14.30 cm(재보정 Stefan), 상수 폭 분할 등각 90\% 구간 ±20.15 cm, 표시 셀의 73.2\% 가 공변량 1개 이상 학습 범위 밖(주로 dem\_rough, dem\_slope, dem\_tpi). 영역 평균 56.09 cm.
  \item 레나델타: 격자 53,011셀(\pth{lena\_grid\_cells\_v1.csv}, 0.009°), 표시 셀 38,086. 블록 교차검증 RMSE 20.77 대 21.48 cm, 90\% 구간 ±25.39 cm, 셀의 84.3\% 가 학습 범위 밖(주로 ERA5-Land 기온·동결 도일). 영역 평균 42.18 cm. 라벨 수별 지도는 라벨 0, 10, 40, 160, 전량 3,037 에서 재보정 계수 1.62 → 2.04(10개), 잔차 가중 1.0 선택(40·160개).
  \item 음영 지형: DEM 을 250 m 로 평균, 방위각 315°, 고도각 45°, 수직 과장 2(음영), 입체 보기는 수직 과장 12(덱 표기; 생성 코드 기본값은 20) [확인: 덱 부제 ‘12배’ 기준으로 말한다].
  \item XM 특징(격자 안 입력): ESA WorldCover 2021 v200 10 m 등급 비율 11열, Sentinel-2 L2A 2019–2023 7–8월 운량 < 30\% 장면에서 NDVI·NDMI 중앙값과 NDVI 표준편차 3열. 라벨 1 km 셀 632개(알래스카 343, 레나 201, 캐나다 88), 타일 89개 × 6장면 = 534장면, 내려받기 5.99 GB, 실패 0. 덮음 WC 100\%, S2 알래스카 100\%, 레나 99.28\%, 캐나다 99.87\%.
  \item AppEEARS(NASA) 대기: XE 2단계의 MODIS 식생(MOD13Q1)·적설(MOD10A1) 점 요청, 레나 육지 요청 78건 처리 중. 마감 10-08 14:22(T0 + 96 h). 새 지역 XF 마감 10-11 14:22(T0 + 7일), 현재 적격 지역 0곳(외부 라벨 2건 Stordalen·Adventdalen 은 북대서양 민감도로만).
\end{itemize}

\section{방법 상세}

\subsection{식}

\begin{itemize}[leftmargin=1.4em, itemsep=0.15em]
  \item Stefan: $\mathrm{ALT} \approx E\,s$, $s=\sqrt{\mathrm{TDD}}$, $E$ 단위 cm·(°C·d)$^{-1/2}$. 물리적으로 $E=\sqrt{2k/(\rho w L)}$($k$ 열전도도, $w$ 수분, $\rho$ 밀도, $L$ 융해 잠열).
  \item 원천 계수: $E_0 = \sum_{i\in S} s_i y_i / \sum_{i \in S} s_i^2$, $S$ = 대상 지역과 100 km 완충 밖의 모든 라벨 셀(원점 통과 최소제곱). 주 대상의 $E_0$ 는 1.59–1.62(알래스카 대상은 1.50, 원천이 레나 중심).
  \item 재보정: $E_n = (n E_{\mathrm{ls}} + \kappa E_0)/(n+\kappa)$, $\kappa = 10$. $E_{\mathrm{ls}}$ 는 대상 라벨 $n$개의 같은 추정량. $\kappa$ 는 앞선 실험에서 정했고 비맹검, 3과 30이 민감도. 현지 최소제곱 Stefan 은 수축 없이 $E_{\mathrm{ls}}$.
  \item 더 강한 기준선: 연도 정합 Stefan = $E_0$ 를 라벨 관측 연도(2010–2024 안)의 √TDD 평균에 적용(라벨 0 에서 원천 계수 대비 −1.00 cm 우세). 라벨 40·160 에서는 토양 물성 Stefan 을 같은 식으로 재보정한 기준선을 병기(잔차 ML 과의 차이 미결정이라 이 $n$에서 ‘재보정 물리식을 넘는다’고 쓰지 않음).
  \item 앵커 구조: $\hat{y}(x) = a(x) + \lambda g(x)$. 물리 잔차 결합은 $a = E_n s$, 잔차 목표는 원천 행 $y - E_0 s$, 대상 행 $y - E_n s$. 원천 행과 대상 라벨 $n$개를 같은 가중으로 학습. $\lambda \in \{0.25, 0.5, 1.0\}$ 을 적합 뒤 곱하며 판정 기준은 잔차 구조 0.25, 직접·물리 입력 구조 1.0. 지역 안 시험과 지도는 라벨 안 블록 교차검증으로 $\lambda$ 선택.
  \item 물리 유사라벨 증강: 대상 라벨 절반(라벨 블록)의 셀에 $E_n s$ 값을 원천 행당 10개 더해 학습. 뽑힌 셀은 실측 유지. 위약 4종은 값만 바꿈(순서 섞기, 대상 풀 평균, 원천 평균, 원천에 적합한 TDD 선형식).
  \item 물리 입력 ML: 직접 ML 입력에 Kudryavtsev 형 모형과 토양 물성 Stefan 출력(또는 앵커 값)을 추가. Stefan·CCI 앵커는 $E_n s$ 와 CCI 값의 평균.
  \item 편향 진단: 추출된 라벨 10개에서 원천 계수 Stefan 의 평균 오차 절댓값 $|\overline{E_0 s - y}|$, 추출·분할 평균. 재보정도 쓰는 라벨만 쓰므로 학습기 적합 전에 안다. 라벨 전량에서 선정 규칙의 원천 계수 대비 이득과의 Spearman ρ 0.38 [0.17, 0.55](28대상), 부호 있는 편향 ρ −0.01.
  \item 이득 분해: 채점 셀을 (0.5° 블록, 기온 √TDD) 묶음으로 나누면 Stefan 예측과 기후 공변량 8종이 묶음 안에서 상수. SSE = Σ(묶음 안 편차 차)² + Σ $N_g$(묶음 평균 차)². 격자 안 RMSE 는 2셀 이상 묶음, 감소 몫은 SSE 로 계산(RMSE 차는 가법적이지 않음). 토양 √TDD 묶음은 민감도(조정 Rand 0.999·0.997·0.832).
\end{itemize}

\subsection{방법 선정 규칙(규칙 W)과 배치 절차}

선정 규칙: (1) 라벨 10개 미만이면 재보정 Stefan. (2) 라벨의 블록을 최대 5겹으로 무작위 배정(블록 1개면 셀 겹), 2겹 미만이면 재보정 Stefan. (3) 겹마다 $E_n$ 재계산 뒤 후보 7개 적합: 원천 계수·재보정·현지 최소제곱 Stefan, 물리 잔차 결합 λ 0.25·1.0, 증강 + 앵커 + 잔차, 물리 유사라벨 증강. (4) 교차검증 RMSE 최소 후보(동점이면 작은 λ, 앞 후보). (5) 전체 라벨로 재적합. 라벨 수 격자 결과를 본 뒤 등록했고 30대상(전이 27 + 새 지역 3)에 적용, 고정 잔차 레시피 대비 비열등 한계 0.5 cm.

사전 지정 배치 절차(Algorithm P): 후보 셀의 표준화 공변량, 블록 크기 $N_b$, 중첩 예산 10 → 40 → 160, 배분 지수 γ. 블록 중심의 최원점 순서로 블록을 돌며 $q_b = n_k N_b^{\gamma}/\sum N_{b'}^{\gamma}$ 개를 배정하고, 블록 안에서는 중심에 가장 가까운 셀(라벨 없을 때) 또는 기존 라벨에서 가장 먼 셀을 고른다. 후보(γ 0·0.5·1, 첫 셀 변형, 블록 층화, 최대 최소 거리)는 알래스카 계열의 개발 분할에서만 기계적 규칙으로 고르게 했고, 규칙에서 후보가 모두 제외되어 Algorithm P = 셀 무작위가 되었다. 따라서 워크플로의 배치 단계 기여는 0 이고 관측 우선순위 지도는 만들지 않았다(등록 규칙).

학습형 배치 정책(DeepSets, 탐색·SI): 알래스카 계열에서 학습해 동결(GPU 결정론 설정, epoch 29, 선택 Jaccard 1.000)한 뒤 알래스카 밖 과제 5개(LE-1, LE-2, CA-2, CA-3, 티베트)에서 한 번 시험. 판정어·p 없이 방향만 적는다. 셀 무작위 대비 상대 RMSE 변화 $v$ 는 LE-1 +0.408, LE-2 −0.027, CA-2 −0.008(가중 불일치), CA-3 −0.036, 티베트 +0.132. 오라클 대비 후회는 CA-2 0.084, CA-3 0.061, LE-2 0.034 로 무작위보다 작고 LE-1 0.344, 티베트 0.195 로 크다. 효용이 섞기 seed 에 민감했다.

\subsection{학습기와 설정}

CatBoost 주 설정 반복 200, 학습률 0.05, 깊이 3, l2 3.0, seed 0·1, CPU, \pth{thread\_count} 명시. 용량 확대 설정 반복 600, 학습률 0.03, 깊이 6. 라벨 0 직접 ML 비교 학습기 10종: CatBoost 3설정(주, 확대, 원천 교차검증 조정), 랜덤 포레스트, TabPFN v2, TabICL v2(미세조정 없음), 신경망 4종(MLP, 다중 헤드 MLP, 축소 FT-Transformer, RealMLP; 초모수는 원천 지역 하나 제외 교차검증). 10종 모두 주 4지역 평균에서 원천 계수 Stefan 보다 열세(+0.85 – +5.21 cm)였고, 두 CI 판정이 일치한 학습기는 5종(RF, 신경망 4종), 나머지는 분할 독립 가정 의존.

\subsection{평가 단위와 통계}

전이 설계: 대상 지역과 100 km 완충 제외. 0.5° 블록을 무작위 순열로 셀 수 균형 절반 A·B 로 5회 분할, A 에서 라벨 추출(seed 고정), B 의 \pth{eval\_mask} 셀에서 채점. 라벨 수 격자 {0, 3, 10, 40, 160, 320, 1000, 전량}, 추출 5회, seed 2개. 러시아 서부·동부는 A 셀 14–17개라 40개 이상 풀은 레나·캐나다. 27대상 = 6지역 + 알래스카 + 하위 지역 10 × 2모드(부모 제외, 부모 유지). 주 4지역 평균이 확인적 풀, 알래스카는 참조(3지역 보조 열), 러시아 중부·티베트는 확장 풀. 지역 안 설계: A 의 라벨만 학습 행, 분할 25회(레나 24 유효), λ 는 라벨 안 블록 교차검증. 첫 시험(분할 5회, 비교 상대 최선 물리 보정, SAR 추가)은 기각되었고 두 번째 시험이 비교 상대·분할·레시피·라벨 범위 네 가지를 바꿨다(둘 다 보고). RMSE 는 셀 가중(주)과 블록 등가중. 두 팔의 라벨 집합이 다른 대비(배치, 워크플로)는 추출도 재표집하는 2단 부트스트랩. 지역 일반 문장은 Hartung–Knapp 구간이 0 을 제외할 때만(지역 4개이면 부호 검정 최소 p 0.125). 오차 하한 = 거친 비지형 공변량이 같은 5셀 이상 묶음 안 관측 분산(알래스카 11.31 cm, 13,333셀; 레나 14.83; 캐나다 17.90).

\section{결과가 왜 이렇게 나왔는가}

\begin{itemize}[leftmargin=1.4em, itemsep=0.25em]
  \item \textbf{라벨 몇 개에서 재보정이 지배하는 이유}: 지역 사이 계수 차이(기하 평균 1.31–11.03)가 Stefan 오차의 가장 큰 성분이고, 계수는 모수 하나라 라벨 10개면 추정된다. 잔차 학습은 공변량과 잔차의 관계를 배워야 하는데 새 지역에서는 원천 관계가 맞지 않고(라벨 10개 이하에서 물리 입력 구조가 잔차 구조보다 2.5–3.7 cm 나쁨), 저가중 λ 0.25 는 보수적이라 몫이 작다. 재보정을 넘어선 ML 몫과 계수 오차의 관계는 확인하지 못했다(사후 분석 ρ 0.28, 주 CI 0 포함).
  \item \textbf{재보정이 캐나다에서 나빠진 이유}: 지역 평균 계수는 원천과 거의 같다($E_0$ 1.59, 지역 자체 1.60). 그러나 캠페인별 계수가 1.06–2.16 으로 라벨 위치마다 달라, 라벨 블록(A)에서 맞춘 계수가 채점 블록(B)에 맞지 않았다(무작위 분할 200회 평균 +3.53 cm). 라벨 단위 민감도에서 캐나다 라벨 10개의 재보정 이득은 0.86 cm 달라졌다. 그래서 라벨 40개부터는 고정 레시피가 아니라 라벨로 고르는 규칙이 필요하다.
  \item \textbf{잔차 ML 이득이 격자 사이에서 나는 이유}: Stefan 예측과 기후 공변량 8종은 ERA5 0.1° 격자 안에서 상수다. 격자 안 변동(알래스카 Stefan 오차 제곱의 72\%)을 맞히려면 지형·토양 공변량이 격자 안 ALT 변동과 상관해야 하는데, 라벨은 1 km 위치 안에서도 표준편차 11–18 cm 로 흩어지고(점 관측, 위치당 행 8.7–39.7), 잔차의 공간 상관 범위는 짧다. 그래서 ML 은 격자 평균의 치우침만 고치고 격자 안 설명 비율은 1\% 미만(알래스카), 8\% 미만(세 지역)이다. 1 km 지도의 분산 97\%(알래스카)가 0.1° 평균으로 설명되는 것도 같은 현상이다.
  \item \textbf{격자 안 입력이 도움이 되지 않는 이유}: 지형 수문 지표·수면·수목 피복 1 km 창 10열은 모든 라벨 수에서 동등, 현장 토양 수분은 재보정 Stefan 의 격자 안 잔차를 설명하지 못했고(블록 교차검증 설명 비율 −4.9 – −3.5\%), WorldCover·Sentinel-2 14열은 격자 안 −0.15 cm(0.5 cm 미만). 1 km 셀 안에서 상수인 특징은 격자 안 분산 가운데 1 km 셀 사이 몫(알래스카 18.8\%, 레나 47.8\%, 캐나다 17.7\%)만 겨냥할 수 있고 나머지는 셀 안 변동이다. 북위 60° 이북에 1 km 미만 토양 수분 위성 제품은 없다. XM 은 남은 변동의 원인을 구별하지 않는다.
  \item \textbf{배치 효과가 지역에 따라 다른 이유}: 캐나다는 블록 39개에 라벨이 퍼져 있고 블록 사이 공변량 변동이 커(블록 사이 비율 0.83) 분산 배치가 새 블록을 덮는다. 레나델타는 유효 블록 1.72개 수준으로 라벨의 79.8\% 가 상위 3블록에 몰려 있고, 블록마다 고르게 뽑으면 소수 셀뿐인 주변 블록이 과대표되어 계수가 치우친다(지역 안 라벨 100개 +1.65 cm, 전이 라벨 10·40개 셀 가중 +1.89·+2.18 cm). 원인 분석은 하지 않았고 이는 결과를 본 뒤의 해석이다. 그래서 ‘어느 지역에서도 나빠지지 않았다(안전)’는 철회되었다.
  \item \textbf{학습된 배치 정책이 옮겨지지 않은 이유}: 알래스카 계열(하위 지역 6개 + 전체)에서만 학습했고 후보 풀은 이미 측정한 셀이다. 알래스카 밖 과제 5개 가운데 두 가중 모두 오차가 작은 과제는 2개(LE-2, CA-3; 계열 기준 0/3), 큰 과제는 2개(LE-1 +0.408, 티베트 +0.132; 계열 2/3)였고 CA-2 는 가중 사이 불일치였다. 효용은 섞기 seed 에 민감했다. 블록 구성이 단순한 알래스카 하위 지역(유효 블록 1.0–1.3)으로는 정책이 배울 구조가 적다. 규칙 기반 Algorithm P 도 알래스카 선택 규칙에서 후보가 모두 제외되었다.
  \item \textbf{CCI v5 가 라벨 셀에서 편향된 이유(단정하지 않음)}: CCI 는 지상 관측이 아니라 모형 화소의 연 최대 융해 깊이이고 재분석으로 구동된다. 라벨 셀에서 편향은 알래스카 +35.63, 캐나다 +58.60 cm(1 km 공통 집합), v4 는 −5.09·+10.35 cm 로 판이 다르면 부호도 다르다. 제품 ALT 정의(모형 화소 최대 융해 깊이 대 GPR·탐침 점 관측)와 기간(1997–2023 평균)이 다르다. 추출 점검(zip md5, 해석식 재구성)은 이상이 없었다. 원인은 ‘제품 성질로 기록’하고 단정하지 않는다.
  \item \textbf{무작위 분할이 직접 ML 을 과대평가하는 이유}: 라벨이 1 km 위치에 몰려(알래스카 위치당 39.7행) 무작위 셀 분할에서는 시험 셀의 최근접 학습 셀 거리가 0 km, 블록 분할 26 km, kNNDM 59 km(알래스카)다. 가까운 라벨을 외우는 직접 ML 은 거리가 멀어질수록 오차가 커지고(11.54 → 17.30 cm), 지역 정보를 쓰지 않는 Stefan 은 거의 변하지 않는다(14.24 → 14.53 cm). 확률 표본이면 무작위 검증도 타당하다는 반론(Wadoux 2021)이 있으나 이 라벨은 확률 표본이 아니다. ‘가장 정직한 추정’이라는 표현은 쓰지 않는다.
  \item \textbf{티베트에서 직접 ML 이 물리식보다 나은 이유}: 원천 계수 Stefan RMSE 244.5 cm 로 계수가 7배 가까이 틀린 심부 레짐(지역 기하 평균 11.03)이고 채점 셀 고도가 원천 99.5 백분위 위라 외삽이다. 계수가 크게 틀리면 재보정(라벨 3개 201 cm, 10개 148 cm)과 잔차(라벨 10개 −21 cm)의 이득도 크다. 독립 지역 근거로는 쓰지 않는다.
\end{itemize}

\section{발표에서 말하지 않은 것}

{\tabsmall
\begin{longtable}{@{}>{\raggedright\arraybackslash}p{3.6cm}>{\raggedright\arraybackslash}p{12.0cm}@{}}
\toprule
항목 & 이유와 현재 상태 \\
\midrule
\endfirsthead
\toprule
항목 & 이유와 현재 상태 \\
\midrule
\endhead
\bottomrule
\endfoot
격자 안 입력 2단계(위성 식생·적설) & AppEEARS 자료 마감 10-08 14:22 까지 대기. 1단계(지형 수문·수면·수목 피복)는 동등, 현장 토양 수분 진단은 설명력 없음. 2단계 표(xt2)는 열지 않았다. 원고에 [PENDING] 3건. \\
새 독립 지역(XF)과 외부 계열 확인 시험 & 마감 10-11 14:22. 공개 자료 조사(DataONE 668건, PANGAEA 166건 등)에서 적격 새 지역 0곳. 외부 라벨 2건은 북대서양 민감도로만. 원고 [PENDING] 3건(XF 2, XC-F3 1). \\
xt2 열람 사건 & 10-05 약 17:00, 덱 그림용 등록 가설 수를 세는 보조 작업에서 봉인 표 \pth{xe\_r1b\_xt2\_tests.csv} 의 판정별 행 수가 화면에 출력되었다(값·CI 는 출력되지 않음). 조치: XE-a·XE-b 에 ‘부분 비맹검(판정 수 열람, 값 미열람)’ 표지, 2단계 집계는 계획대로, 보조 작업이 봉인 폴더를 읽지 않도록 지시문 수정(커밋 2047e77). 질문이 있으면 그대로 설명한다. \\
비용 & 프로젝트 누적 약 225 달러(장부 정산 220.26 + WF 약 2.98 + 작업 A 약 1.40). 과학적 결론과 무관하므로 말하지 않는다. \\
작업 방식 & 실험 실행·문서 작성의 작업 도구는 과학적 내용과 무관하므로 발표와 질의응답에서 언급하지 않는다. \\
관측 우선순위 지도 & 등록 규칙(WRAPUP 11.1·11.2)에 따라 만들지 않았다. 배치 알고리즘이 무작위로 귀결되었고 음성 결과가 누적되어 있다. \\
신경망 학습기 제외 & 환경 간 최대 5.6 cm 차이로 최종 묶음에서 제외. 라벨 0 직접 ML 비교에만 포함. \\
다중 원천 적층(XB) & 전이 가설 2개 기각(미결정), 지역 안 라벨 500개 −0.55 cm 부분 지지, 알래스카 비열등 지지. CCI v5 후보를 더한 민감도에서 부분 지지가 유지되지 않음. 덱 부록 18. \\
기후 외삽 재시험(XI) & 두 가설 기각, 한 가설(외삽 블록에서 잔차 결합이 재보정 Stefan 보다 −2.44 cm) 지지. 같은 시기의 따뜻한 블록을 쓴 공간 대용이라 시간 외삽 근거로 쓰지 않음. \\
지온 유도 보조 라벨(XJ) & 티베트 6행 모두 우세(−8 – −83 cm)이나 정의가 직접 관측과 달라 일반화하지 않음. 원천 쪽은 약관 근거 미기록으로 시험하지 않음. \\
대회 결과 & 예선 통과·본선 순위·수상 기록이 저장소에 없다. 묻는다면 ‘기록 없음’으로 답한다. \\
\end{longtable}}

\section{예상 질문과 답}

\begin{enumerate}[leftmargin=1.8em, itemsep=0.45em, label=\textbf{Q\arabic*.}]
  \item \textbf{결국 Stefan 식에 보정을 붙인 것 아닌가.} 그렇다. 그것이 결론이다. 라벨 0개에서는 원천 계수 Stefan 이 직접 ML(학습기 10종 모두 열세)과 기존 지도(CCI v5 +23.8 cm)보다 낫고, 라벨 10개의 이득 2.45 cm 는 계수 하나를 고친 것이며, 잔차 ML 의 순가치는 라벨 전량에서 0.40 cm(주 4지역), 알래스카 지역 안 0.54 cm 다. 기여는 이 사실을 라벨 수 축 위에서 같은 시험지로 보인 평가 설계와 그 결과로 만든 워크플로이지 새 모형이 아니다.
  \item \textbf{학습 원천이 78–94\% 알래스카인데 일반화할 수 있나.} 일반화 범위를 그만큼으로 한정한다. 주 4지역 판정은 재사용 지역의 재검정이고 확인적 독립 지역은 4–5개(러시아 중부 확충판 포함)다. 지역 홀드아웃으로 외삽을 평가했고 지역 일반 문장은 Hartung–Knapp 구간이 0 을 제외할 때만 쓴다. 새 지역 확보(XF, KOPRI·KPDC 협업)가 다음 단계다.
  \item \textbf{0.5 cm 가 의미 있는 차이인가.} ±0.5 cm 는 측정 오차가 아니라 결과 전에 고정한 동등 한계이며 주 4지역 원천 Stefan RMSE(21.7–43.0 cm)의 약 1–2\% 다. 효과 크기는 오차 하한과 함께 보고한다. 알래스카 라벨 1,000개의 0.54 cm 는 RMSE 3.7\% 이지만 하한 11.31 cm 위의 줄일 수 있는 오차 제곱으로는 19\% 감소다. 민감도 한계 1.0 cm 와 RMSE 의 2\% 판정도 함께 냈다.
  \item \textbf{왜 딥러닝을 쓰지 않았나.} 썼다. 라벨 0 직접 ML 에 신경망 4종(원천 지역 하나 제외 교차검증으로 조정)과 표형 파운데이션 모델 2종을 넣었고 모두 원천 계수 Stefan 보다 열세였다(+0.85 – +5.21 cm). 또 신경망은 계산 환경 사이에서 최대 5.6 cm 달라 재현 관문을 통과하지 못했으므로 최종 묶음은 CatBoost 와 닫힌 물리식만 쓴다. 자료가 약 1.7만 행 × 25열이라 CatBoost 한 번이 수 초이며 적합 최소 86만 건의 대부분을 CPU 로 돌렸다(파운데이션 모델·신경망은 로컬 GPU).
  \item \textbf{극지연구소 자료가 들어오면 무엇이 달라지나.} 두 가지다. 약관이 확인된 새 독립 지역이 생기면 라벨 10개 워크플로의 ‘실행 전 등록 확인 시험’(XC-F3)과 핵심 대비 재계산(XF)을 그대로 적용할 수 있다. 지중온도 프로파일이 있으면 지온 유도 라벨의 정의 차이(콘슬에서 탐침 35.0, 코어 81.3, 유도 136.4 cm 로 갈림)를 같은 지점에서 비교할 수 있다.
  \item \textbf{정보 해상도가 약 10 km 라는 주장의 근거는.} 셋이다. (1) 이득 분해: 알래스카 잔차 ML 의 오차 제곱 감소 99\% 가 ERA5 격자 사이 성분, 격자 안 설명 1\% 미만. (2) 표시 해상도 시연: 1 km 지도 분산의 97\%(알래스카)·61\%(레나)가 0.1° 셀 평균으로 설명, 0.1° 셀 안 1 km 표준편차 1.8·1.1 cm. (3) 격자 안 입력(10 m–1 km) 추가가 격자 안 오차를 0.5 cm 이상 줄이지 못함. 그래서 ‘1 km 셀로 표시한 기후 격자 단위 보정 지도’라고 쓴다.
  \item \textbf{왜 사전 등록을 했나, 내부 등록이 효력이 있나.} 7월 대회 때 같은 교차검증에서 고른 최선 설정(13.33 cm, 22.92 cm)이 재평가에서 사라지는 것을 겪었다. 그래서 가설·판정 규칙·허용 문장을 결과 전에 커밋하고 봉인·열람 순서·재현 관문을 두었다. 제3자 타임스탬프가 없는 내부 등록이라 원고에 그렇게 밝히고 커밋 해시와 시각을 SI 에 적는다. 초록 판정어는 사전 등록 대비 10개에만 쓴다.
  \item \textbf{κ = 10 과 λ 는 어떻게 정했나.} κ 는 이 연구의 앞선 실험에서 정한 값이고 비맹검임을 원고에 밝혔다. 라벨 몇 개의 재보정이 오차를 키울 수 있다는 근거(Steyerberg 2004 류)로 원천 계수에 라벨 10개 무게를 준다. 3과 30 은 민감도. κ 를 오프셋 최대우도, 혼합효과 부스팅, PPI++ 로 대체하는 부분 풀링 시험(09-26)은 기각되었고 14대상 평균에서 κ = 10 수축이 최선이었다(라벨 3개 −0.62, 10개 −1.16 cm). λ 는 {0.25, 0.5, 1.0} 을 적합 뒤 곱하고, 전이 판정 기준 0.25 는 본 실행 전에 고정했으며, 지역 안과 지도에서는 라벨 안 블록 교차검증으로 고른다.
  \item \textbf{워크플로가 독립적으로 검증되었나.} 아니다. 워크플로 순차 적용(XC)은 구성 요소를 설계한 지역을 새 seed(201–210)로 다시 나눈 재검정이며 ‘재사용 지역의 재검정’ 표지를 단다. 배치 단계는 무작위로 귀결되어 기여가 없고, 160개 이득은 캐나다에서 났으며 알래스카·레나·러시아 동부는 일부 라벨 수에서 0.5 cm 넘게 나빴다. 독립 검증은 새 지역(XF) 확보 뒤의 일이다.
  \item \textbf{캐나다에서 재보정이 나빠진 것은 방법 결함 아닌가.} 수축 재보정 자체의 결함이라기보다 라벨 위치 이질성의 문제다. 캠페인별 계수가 1.06–2.16 으로 다르고 지역 평균 계수는 원천과 같다. 현지 최소제곱(수축 없음)은 더 나쁘다(레나 라벨 3개 24.8 대 21.8 cm). 그래서 40개부터 라벨로 방법을 고르게 했고, 캐나다에서는 선정 규칙이 고정 레시피보다 2.1–2.7 cm 낫다.
  \item \textbf{물리 유사라벨 증강은 결국 효과가 없는 것인가.} 라벨 0개 전용이다. 순서를 섞은 유사라벨보다 1.63 cm 낫고(사전 등록, Holm p 0.001) 직접 ML 의 큰 실패(2 cm 초과 악화 18/30)를 2/30 으로 줄인다. 다만 선형 융해 지수 위약과의 차이는 0.48 cm 이고 라벨 10개부터는 재보정이 같은 몫을 가져가 추가 효과가 없다(증강 + 앵커 + 잔차 − 잔차 −0.13 cm 동등).
  \item \textbf{기존 지도와 비교가 공정한가.} 같은 채점 셀, 같은 라벨, 학습 지점 5·25 km 블록 마스크 뒤 비교이고 제품 ALT 정의가 다르다는 점을 함께 적는다. Wei 2026 은 대상 지역 안 CALM 지점을 학습에 썼으므로 미결정 결과만 적고 우열을 말하지 않는다. CCI v5 는 학습 지점이 없어 마스크에 무관하며 라벨 0 에서 +23.79 cm 열세(Holm p 0.0006), 라벨 10개를 같은 라벨로 재보정해도 잔차 결합이 8.37 cm 낫다. 등록 이탈(WRAPUP 10) 표지를 단 분석이다.
  \item \textbf{예측 구간은 믿을 수 있나.} 알래스카 지역 안에서는 포함률 0.45 → 0.93 으로 보정되지만 셀별 폭의 구간 점수(68.85)는 상수 폭(65.39)보다 낫지 않다. 새 지역에 알래스카 보정 구간을 그대로 쓰면 0.32–0.73 으로 떨어지고 계층 conformal 은 0.86 [0.80, 0.92]이나 집단이 4개라 유한표본 보장이 없다. 지도의 90\% 구간은 상수 폭(알래스카 ±20.2 cm)이며 외삽 셀(73\%)에서 포함률을 보장하지 않는다.
  \item \textbf{라벨 수 n 의 뜻이 지역마다 다른데 곡선을 비교해도 되나.} 그래서 n 을 ‘행 수’로 정의하고 행의 뜻(점 관측, CALM 지점 다년 평균, 1 km 셀 평균)을 지역마다 밝힌다. 라벨 단위 민감도(1 km 위치 평균 대 점)에서 캐나다 라벨 10개의 재보정 이득이 0.86 cm 달라졌고, 단년 대 다년 라벨 차이는 0.05 cm 이하였다. 지역 평균 판정을 지역 단위 판정처럼 쓰지 않는다.
  \item \textbf{무작위 분할이 맞는 평가 아닌가.} 질문에 따라 다르다. 학습 지역 안 보간이 목적이면 무작위도 타당하고 확률 표본이면 더욱 그렇다(Wadoux 2021). 이 연구의 질문은 새 지역이고 라벨은 관측망에 몰린 비확률 표본이므로 지역 홀드아웃과 kNNDM 을 쓴다. 무작위 수치(직접 ML 14.8 cm)도 함께 싣는다.
  \item \textbf{알래스카 결과가 대회 결과와 모순되지 않나.} 모순이 아니다. 대회의 7.8\%(14.46 → 13.33 cm)는 같은 교차검증에서 고른 최선 설정이고, 이번 3.7\%(14.41 → 13.87 cm)는 설정을 라벨 안에서 고른 뒤 블록 밖에서 채점한 값이다. 줄일 수 있는 오차 제곱으로는 39\% 대 19\% 로 같은 방향이다. 알래스카를 새 지역으로 두면(원천 약 3,860셀) 라벨 전량에서도 차이를 확인하지 못했다.
  \item \textbf{그림에 왜 판정 기호나 p 값이 없나.} 판정은 두 가중 신뢰구간과 ±0.5 cm 띠로 읽도록 했고 Holm 보정 p 는 SI 표 S1 에 모았다. 초록 판정어는 사전 등록 대비 10개(방향 확정 6, 동등 1, 보정 전만 유의 1, 미결정 2)에만 쓴다.
  \item \textbf{라벨이 많은 캐나다에서는 직접 ML 이 우리 방법보다 나아 보이는데, 워크플로는 이를 어떻게 다루나.} 캐나다에 한정된 결과다. 캐나다 지역 안(분할 25회)에서 직접 ML − 재보정 Stefan 은 라벨 200개 −1.50 [−1.99, −0.24](블록 등가중 −0.03 [−1.03, 1.01]), 전량 −1.50 [−2.19, −0.12](−0.30 [−1.44, 0.86])로 셀 가중만 0 을 벗어난 미결정이고, 같은 조건에서 직접 ML 도 원천 계수 Stefan 보다는 2.37·2.47 cm 나빴다(재보정 Stefan 은 +3.86·+3.98). 즉 캐나다는 라벨 위치 이질성 때문에 재보정 자체가 손해인 지역이며 알래스카·레나에서는 직접 ML 이 열세다. 워크플로는 라벨 40개 이상에서 대상 라벨 교차검증으로 후보 가운데 오차가 가장 작은 방법을 고르므로 이런 지역 차이를 자동으로 반영하고, 캐나다에서 선정 규칙은 고정 레시피보다 2.10·2.72 cm 나았다. 다만 현재 후보 7개에 직접 ML 은 없다. 후보에 직접 ML 을 더하는 것은 다음 단계이며, 라벨 0–10개의 열세를 피하기 위해 40개 이상에서만 후보로 두어야 한다.
  \item \textbf{라벨 10개 재보정이 캐나다에서는 왜 손해인가.} 캐나다 라벨 10개 재보정 − 원천 계수 Stefan 은 +2.26 [0.71, 2.94] cm 로 오차 증가다. 지역 평균 계수는 원천과 거의 같지만(원천 1.59, 지역 자체 1.60) 캠페인별 계수가 1.06–2.16 으로 라벨 위치마다 달라, 라벨 블록에서 뽑힌 10개로 맞춘 계수가 채점 블록에 맞지 않는다(무작위 분할 200회 평균 +3.53 cm, 라벨 단위 민감도 0.86 cm). 수축을 없앤 현지 최소제곱은 더 나쁘다. 워크플로의 라벨 10개 편향 진단은 원천 계수 Stefan 의 평균 오차 절댓값이며 라벨 기반 보정 이득의 크기와 ρ 0.38 로 상관하므로, 캐나다처럼 지역 평균 계수가 맞는 곳에서는 작은 진단값이 ‘재보정으로 얻을 이득이 작다’는 신호가 된다(부호는 예측하지 않음 ρ −0.01, 몫별 관계는 확인되지 않음). 라벨 40개부터는 선정 규칙이 캐나다에서 다른 후보를 골라 손해를 되돌리며, 순차 적용 결과의 캐나다 행은 재보정 Stefan 대비 라벨 10개 −0.84, 40개 −2.96, 160개 −3.55 cm 였다.
  \item \textbf{다음에 무엇을 하면 결과가 바뀔 수 있나.} 격자 안 변동을 설명할 입력(토양 수분·유기층·미지형 30–500 m)과 새 독립 지역이다. 공개 10–500 m 입력은 이득을 더하지 못했고 북위 60° 이북에 1 km 미만 토양 수분 제품은 없다. 위성 적설·식생(10-08)과 새 지역(10-11)이 들어오면 같은 하네스로 재계산한다.
\end{enumerate}

\clearpage
\section{수치 요약표}

{\tabsmall
\begin{longtable}{@{}>{\raggedright\arraybackslash}p{1.2cm}>{\raggedright\arraybackslash}p{5.7cm}>{\raggedright\arraybackslash}p{2.9cm}>{\raggedright\arraybackslash}p{5.0cm}@{}}
\toprule
쪽 & 수치 & 종류 & 출처 \\
\midrule
\endfirsthead
\toprule
쪽 & 수치 & 종류 & 출처 \\
\midrule
\endhead
\bottomrule
\endfoot
1, 7 & 직접 라벨 1 km 셀 17,467(표 17,572행), 원천의 알래스카 비율 78–94\% & 서술 & D README E04, E35 \\
1, 33 & 알래스카 영역 평균: 우리 56.1, 재보정 Stefan 57.5, CCI v5 105.5, Wei 70.8, Aalto 78.8, Yi·Kimball 105.1 cm; r 0.96 / 0.46 / −0.15 / −0.10 / 0.34 & 서술 & XK/XL 추가 등록 결과 표(xl\_summary\_v1.csv) \\
1, 9 & 지역 계수 기하 평균 1.31(레나)–11.03(티베트) & 서술 & D README E38 \\
7 & 지역별 행/위치/블록: 레나 3,037/201/20, 캐나다 750/86/39, 러시아 W 31/29/21, E 30/30/21, 알래스카 13,606/343/74, 러시아 C 57/57/13, 티베트 132/132/34 & 서술 & Table 1, D README E10–E18 \\
9 & 러시아 서부 예시 셀 TDD 1,537 °C·d, 실측 118 cm, 원천 예측 62 cm; 원천 1.59 → 재보정 2.11 & 예시 & 덱 9쪽 그림 표기 \\
10 & 레나 예시: 최소제곱 1.352, 재보정 1.486, 원천 1.619 & 예시 & 덱 스펙 N07 참고 줄 \\
11 & CatBoost 200회·학습률 0.05·깊이 3·seed 2; λ {0.25, 0.5, 1.0}, 기준 0.25 & 설정 & methods.tex, h40 cb\_fit \\
13 & 분할 5 × 추출 5 × seed 2, 라벨 격자 8단계, 재표집 10,000(격자 집계 1,000) & 설계 & D README E24–E25, methods.tex \\
14 & 직접 ML 3지역 평균 14.81(무작위) → 28.71 cm(지역 홀드아웃); 재보정 Stefan 20.54–22.24; 무작위→kNNDM 증가 직접 5.76/8.17/14.20, Stefan 0.29/0.51/4.28 & 서술 & Fig1\_values.txt, 계획서 8.3 \\
16, 18 & 라벨 0에서 2 cm 초과 악화 대상: 직접 18/30, 증강 2/30(CI 기준 14, 8); 알래스카 직접 35.33, 증강 16.33, 원천 Stefan 14.54 cm; 증강 − 원천 +1.79 & 사후 서술 & C1 README 2.5–2.8, wf0\_risk.csv \\
17 & 물리 유사라벨 − 섞은 유사라벨 −1.63 [−1.98, −1.04], 블록 등가중 −1.05, Holm p 0.001; 선형 융해 지수 대비 −0.48(라벨 0), −0.08 동등(라벨 10) & 사전 등록 & C1 README 2.3(AB3, L15) \\
18 & 티베트 원천 Stefan RMSE 244.51 cm; 직접 ML − 원천 −59.98 cm & 서술(교차 환경, 비맹검) & SI\_learners README, C2 README TB6 \\
19 & 연도 정합 Stefan − 원천 −1.00 [−1.35, −0.32]; 직접 ML − 연도 정합 +3.24 [1.86, 4.29]; CCI v5 − 원천 +23.79 [15.34, 30.16] / +19.17, Holm p 0.0006; 라벨 10 잔차 − 재보정 CCI −8.37 & 등록(CCI 는 등록 이탈 표지) & C1 README 2.7, 계획서 8.7 \\
20 & 재보정 − 원천, 라벨 10, 주 4지역 −2.45 [−3.04, −1.88], 블록 −2.62, Holm p 0.001; 러시아 W −11.96; 캐나다 +2.26 [0.71, 2.94]; 레나·캐나다 풀 라벨 3·10 +0.56, +1.24 & 사전 등록 & C2 README A6, CA1; pool\_fixed\_curve.csv \\
21 & 잔차 결합 − 재보정, 라벨 10 −0.18 [−0.53, −0.01], p 0.034, Holm 0.136; 전량 −0.40 [−0.76, −0.22]; TabPFN 잔차 −0.32(3개), TabICL −0.55(10개) & 사전 등록 & C2 README A7·A9, C3 README \\
22 & 선정 규칙 − 고정 레시피: 40개 −0.87 [−1.20, −0.27](분할 독립 의존), 160개 −1.32 [−1.70, −0.68]; 캐나다 −2.10, −2.72; 레나 +0.35, +0.07; 10개 비열등 불성립(상한 0.73/0.69) & 사후 분석 & C5 README 1.2–1.3, wf\_tests.csv WF4-a \\
23 & 워크플로 − 무작위+고정: 40개 −0.86 [−1.03, −0.39] 미결정, 160개 −0.93 [−1.25, −0.40] Holm 0.040(캐나다 −2.41); 워크플로 − 재보정: 10개 +0.08 [−0.07, 0.20], 40개 −1.35, 160개 −1.51; 0.5 cm 초과 악화 지역: 러시아 E(10개 +1.34), 알래스카(40개 +1.51, 160개 +2.34), 레나(160개 +0.53) & 사후 분석, 재사용 지역 재검정 & 계획서 8.10(xc\_hyp.csv) \\
24 & 알래스카 라벨 1000: 14.41 → 13.87, −0.54 [−0.69, −0.32], 블록 −0.60; 500개 −0.39, 전량 −0.52(공통 CI 미결정); 하한 알래스카 11.31, 레나 14.83, 캐나다 17.90 & 사후 분석 & C4 README, lgx\_floor.csv \\
24 & 줄일 수 있는 오차 제곱 감소 19.2\%(이번) 대 38.7\%(대회) & 서술 & RESEARCH\_OVERVIEW C4 \\
25 & 총 −0.52, 격자 사이 −1.03 [−1.30, −0.59], 격자 안 −0.01 [−0.03, 0.04]; 격자 사이 몫 약 99\%; Stefan 오차 제곱의 격자 안 비율 72\%; 격자 안 설명 알래스카 < 1\%(최대 0.7), 세 지역 < 8\%(최대 7.8); 3지역 평균 격자 안 +0.09 & 사후 분석 & C8 README 2.1–2.3 \\
26 & 캐나다 분산 배치 − 무작위: 50–100개 −2.54 – −3.21, 200개 −1.25 – −1.44, 전이 40개 −2.39 [−3.40, −0.49]; 레나 블록 층화 100개 +1.65 [0.90, 2.40]; 능동 선택 +0.12 – +2.79 & 사후 분석 & C6 README 1.3, Fig5\_values.txt \\
27 & 5지역 라벨 0–40 직접 ML 유의 개선 0/5; 전량 잔차 − 원천 −3.19 [−4.19, −2.22]; 라벨 10 잔차 − 재보정 −0.13 동등; 러시아 C 전량 −5.42 [−8.96, −1.64] / −2.55 [−6.83, 1.48] 미결정 & 등록(교차 환경) & SI\_learners README 1.3 \\
28 & 계층 conformal 0.86 [0.80, 0.92](러시아 W 0.75); 알래스카 보정 구간 전이 0.32–0.73; 생성 구간 0.76, 0.65; 알래스카 CQR 0.45 → 0.93, 구간 점수 68.85 대 65.39 & 보조 & SI\_uncertainty README \\
29 & 표시 셀 858,517; CV RMSE 13.80 대 14.30; 90\% 구간 ±20.15; 범위 밖 셀 73.2\%; 재보정 계수 1.6166, λ 0.5 & 서술 & alaska\_pred\_v1\_meta.json, alaska\_cv\_oof\_v1.csv \\
29, 36 & 알래스카 블록 교차검증: 재보정 Stefan 14.30 → 물리 잔차 결합 13.80 cm, 74개 블록 중 47개에서 잔차 ML 오차가 더 작음 & 서술 & 덱 스펙 NOP(그림 설명문), alaska\_cv\_oof\_v1.csv \\
31 & 레나 라벨 수별: 계수 1.62 → 2.04(10개), λ 1.0(40·160개), 전량 3,037 & 서술 & 덱 스펙 NLS, label\_sequence\_v1 \\
32 & 0.1° 평균이 1 km 분산의 97\%(알래스카), 61\%(레나) 설명; 0.1° 안 표준편차 1.8(전체 9.5), 1.1(전체 1.8) cm & 서술 & XK/XL 추가 등록 결과 \\
34 & 레나 입체: 수직 과장 12, 구릉 최고 약 580 m, 평원 약 30 m, 범위 약 246 × 242 km, 색 35–50 cm & 서술 & 덱 스펙 NL3 \\
35 & 레나델타 라벨 0개, 라벨 셀 3,037개 RMSE: 원천 계수 Stefan 21.7, 직접 ML 25.5(+3.8), 물리 유사라벨 증강 21.9 cm(Stefan 수준) & 서술 & 덱 스펙 NTM(그림 설명문, 등록 기록 재현) \\
37 & 결과 종합: 라벨 0–10 직접 ML − 원천 +2.2 → +1.3, 재보정 −2.45(10개), 잔차 결합은 재보정과 0.5 cm 안; 지역 안 알래스카 −0.54(1000개), 레나·캐나다 +0.2 – +0.6 미결정, 알래스카 직접 ML +2 – +3 & 서술 & 덱 스펙 NR2 참고 줄 \\
38 & 평가 지역 7, 전이 대상 27, 라벨 8단계, 방법 12·계열 5, 학습기 10, 가설 174(문서 7종: 45+11+17+15+24+58+4), 적합 ≥ 863,229, 재표집 10,000; 주 가설 10: 방향 확정 6(감소 5, 증가 1), 동등 1, 미결정 3(덱 표기; E1 원천값은 보정 전만 유의 1 + 미결정 2); 환경 간 차 3.6e-14 cm(34,360키) & 서술 & E1 원천값 JSON(deck/assets/paper\_report/ method/) \\
39 & 강점: 2 cm 초과 악화 18/30 대 2/30; 워크플로 − 재보정 −1.35(40개)·−1.51(160개); 알래스카 잔차 결합 −0.54; CCI v5 대비 라벨 0 −23.8, 라벨 10 −8.4(같은 라벨로 재보정한 CCI v5 대비 잔차 결합 −8.37) & 사후 서술·사후 분석·등록 혼합 & 덱 스펙 NS1, 계획서 8.7(XG-2c) \\
40 & 격자 안 입력(XM): 격자 안 −0.15 [−0.25, −0.05] 우세(Holm 0.019), 1,000개 −0.21, 500개 동등; 총 RMSE 3지역 −0.70, 레나 −1.50, 알래스카 −0.12, 캐나다 −0.49 미결정; 레나 격자 사이 −1.80 & 탐색(SI) & XM 추가 등록 결과 \\
부록 15 & 격자 크기별 재보정 − 잔차(알래스카): 1 km 0.50 [0.18, 1.10], 0.05° 1.67, 0.1° 1.54, 0.25° 1.61; 캐나다 1 km 2.62; 레나 CI 0 포함 & 서술 & XK 추가 등록 결과 \\
부록 18 & 적층(XB): 전이 미결정(기각), 지역 안 500개 −0.55 / −0.33 우세, 1,000개 −0.498, 알래스카 비열등 4개 n 모두 & 결과 열람 뒤 설계 & 계획서 8.2, 8.12 \\
부록 19 & 기후 외삽(XI): XI-c −2.44(100개), −2.41(전량) 지지; 지온 보조(XJ): 티베트 −8.14 – −83.17 & 결과 열람 뒤 설계 & 계획서 8.4, 8.5 \\
\end{longtable}}

\section{출처를 찾지 못해 뺀 수치}

\begin{itemize}[leftmargin=1.4em, itemsep=0.15em]
  \item 대회 예선 통과 여부·본선 순위·수상: 저장소와 기록에 없다. 대본과 배경 모두에서 뺐다.
  \item Rescale 누적 지출의 확정 합계: 장부에 WF·작업 A 7건의 실제 비용이 비어 있어 약 225 달러는 문서 합산값(220.26 + 2.98 + 1.40 ≈ 224.64)이며 이를 적은 문서는 없다. 대본에서는 비용을 말하지 않는다.
  \item 라벨 3개의 주 4지역 재보정 이득 값: 그림(덱 20쪽, 약 −0.9 cm)에만 있고 노트·판정 기록에 수치가 없어 대본에서는 ‘라벨 3개부터 0선 아래’로만 말하도록 했다가 최종판에서는 뺐다.
  \item 덱 24쪽 알래스카 직접 ML 의 라벨 1000개 값(그림상 약 +2.8 cm): 노트에 없어 ‘0선 위’로만 말한다.
  \item 레나 입체 보기의 수직 과장: 덱 부제와 그림 표기는 12배, 생성 코드 \pth{lena\_3d\_v1.py} 의 기본값은 20이다. 대본은 덱 표기 12배를 따른다. 질문이 나오면 코드 기본값과 다를 수 있다고 답한다.
  \item 레나 블록 층화의 ‘+1.7 cm’: 덱 표기이며 원천 값은 라벨 100개 +1.65 cm(두 가중 열세), 20–500개 +1.47 – +1.98 cm 다.
  \item KPDC AWS 자료의 식별자 귀속: 원고 초안에 [확인 필요]로 남아 있어 배경에서 식별자 3개(코어·프로파일)만 적었다.
  \item 덱 38쪽 ‘주 가설 10개’ 행의 ‘미결정 3’은 E1 원천값(동등 1, 보정 전만 유의 1, 미결정 2)의 뒤 둘을 합친 표기다. 대본은 덱 표기를 따르고 선택 상자에 ‘동등 1개, 미결정 3개’로 적었다.
  \item 로컬 scikit-learn 판: \pth{environment-lock.txt} 에 판 번호가 없고 설치판은 1.3.0 이다(추가 묶음 고정판 1.9.1 과 다름). 표에는 둘을 구분해 적었다.
\end{itemize}


\end{document}
